\documentclass[10pt,a4paper]{article}
\usepackage[utf8]{inputenc}
\usepackage[T1]{fontenc}
\usepackage{lmodern}
\usepackage[margin=1.8cm, top=2.2cm, bottom=2.2cm]{geometry}
\usepackage{xcolor}
\usepackage{titlesec}
\usepackage{fancyhdr}
\usepackage{hyperref}
\usepackage{booktabs}
\usepackage{longtable}
\usepackage{array}
\usepackage{enumitem}
\usepackage{needspace}
\usepackage[most]{tcolorbox}
\tcbuselibrary{skins,breakable}

% Color Definitions
\definecolor{sihnavy}{HTML}{0B2545}
\definecolor{sihblue}{HTML}{134074}
\definecolor{softgray}{HTML}{F8FAFC}
\definecolor{bordergray}{HTML}{CBD5E1}
\definecolor{diffextreme}{HTML}{991B1B}
\definecolor{diffveryhard}{HTML}{C2410C}
\definecolor{diffhard}{HTML}{D97706}
\definecolor{diffmod}{HTML}{047857}

\hypersetup{
    colorlinks=true,
    linkcolor=sihblue,
    urlcolor=sihblue,
    citecolor=sihblue,
    pdftitle={SIH 2026 - Software and Disaster Management Problem Statements},
    pdfauthor={SIH 2026 Extractor}
}

\pagestyle{fancy}
\fancyhf{}
\fancyhead[L]{\small\color{sihblue}\textbf{SIH 2026: Software $\vert$ Disaster Management}}
\fancyhead[R]{\small\color{gray}Difficulty Ranked Compendium}
\fancyfoot[L]{\small\color{gray}Source: sih.gov.in/sih2026PS}
\fancyfoot[R]{\small\color{sihblue}\textbf{Page \thepage}}
\renewcommand{\headrulewidth}{0.4pt}
\renewcommand{\footrulewidth}{0.4pt}

\titleformat{\section}{\Large\bfseries\color{sihnavy}}{\thesection}{1em}{}[\titlerule]
\titleformat{\subsection}{\large\bfseries\color{sihblue}}{\thesubsection}{1em}{}

\tcbset{
    pscard/.style={
        enhanced,
        breakable,
        colback=softgray,
        colframe=bordergray,
        arc=2.5mm,
        boxrule=0.7pt,
        left=9pt,
        right=9pt,
        top=9pt,
        bottom=9pt,
        fonttitle=\bfseries,
        coltitle=white,
        colbacktitle=sihnavy,
        attach boxed title to top left={yshift=-2mm, xshift=4mm},
        boxed title style={arc=1.5mm, boxrule=0.5pt}
    }
}

\begin{document}

\begin{center}
    \vspace*{0.4cm}
    {\Huge \textbf{\color{sihnavy}Smart India Hackathon 2026}}\\[0.35cm]
    {\LARGE \textbf{\color{sihblue}Category: Software \quad $\vert$ \quad Theme: Disaster Management}}\\[0.25cm]
    {\large \color{gray} Ranked by Technical Difficulty \& Complexity Score (Out of 10.0)}\\[0.4cm]
    \rule{0.85\textwidth}{1pt}\\[0.35cm]
    {\large \textbf{Total Statements: } 24 \quad $\vert$ \quad \textbf{Difficulty Range: } 4.0 - 8.8 / 10.0}\\[0.15cm]
    {\small \color{gray} Generated on: \today}
\end{center}

\vspace{0.4cm}
\hrule
\vspace{0.4cm}

\section*{Difficulty Ranking \& Overview}
\addcontentsline{toc}{section}{Difficulty Ranking \& Overview}

\begin{tabular}{rlllp{8.5cm}}
\toprule
\textbf{\#} & \textbf{Code} & \textbf{Difficulty} & \textbf{Level} & \textbf{Title} \\
\midrule

1 & \textbf{SIH26084} & \textbf{9.2/10} & Extreme & Convective scale nowcasting for Thunderstorms, Hail \& Cloudbursts... \\

2 & \textbf{SIH26085} & \textbf{9.2/10} & Extreme & Urban Flood Nowcasting System (Drainage and Rainfall Coupling)... \\

3 & \textbf{SIH26077} & \textbf{8.6/10} & Extreme & AI-Driven Hyper-Local Early Warning System for Severe Weather Now... \\

4 & \textbf{SIH26161} & \textbf{8.6/10} & Extreme & Dam Break Inundation Modelling Using Hydrodynamic Modelling of an... \\

5 & \textbf{SIH26072} & \textbf{8.3/10} & Very & AIML based Nowcasting of thunderstorm and lightning using atmosph... \\

6 & \textbf{SIH26001} & \textbf{8.2/10} & Very & AI-Based early warning and landslide Risk Monitoring System in NE... \\

7 & \textbf{SIH26071} & \textbf{8.0/10} & Very & AI/ML-Based Integrated heavy rainfall Early Warning and Inundatio... \\

8 & \textbf{SIH26078} & \textbf{8.0/10} & Very & AI-Driven Spatio-Temporal Tracking of Extreme Weather Anomalies i... \\

9 & \textbf{SIH26083} & \textbf{7.9/10} & Very & Extreme Heatwave Early Warning and Human Thermal Stress Index... \\

10 & \textbf{SIH26013} & \textbf{7.3/10} & Hard & Automated lntegration and lntelligent Harmonization of Multi-sour... \\

11 & \textbf{SIH26028} & \textbf{7.3/10} & Hard & Dynamic Forecast of Expected Time of Arrival (ETA) for Coaching T... \\

12 & \textbf{SIH26073} & \textbf{7.2/10} & Hard & AI/ML-Based Intelligent Anomaly Detection for Automatic Weather S... \\

13 & \textbf{SIH26192} & \textbf{7.2/10} & Hard & Flash Flood Prediction System for Hilly Regions using Multi-Sourc... \\

14 & \textbf{SIH26068} & \textbf{7.1/10} & Hard & WeatherGPT: Conversational AI for Weather Forecasting, Alerts, an... \\

15 & \textbf{SIH26079} & \textbf{7.1/10} & Hard & AI-Based Forecast Bust Detection for Medium-Range Weather Forecas... \\

16 & \textbf{SIH26080} & \textbf{7.1/10} & Hard & Regime-Aware AI Post-Processing of Monsoon Rainfall Forecasts... \\

17 & \textbf{SIH26082} & \textbf{7.1/10} & Hard & Air PollutionWeather Coupled Forecasting System (Delhi NCR Focus)... \\

18 & \textbf{SIH26015} & \textbf{6.5/10} & Hard & Application of Geospatial Techniques for visualization and analys... \\

19 & \textbf{SIH26043} & \textbf{6.5/10} & Hard & A digital platform to crowdsource societal challenges and facilit... \\

20 & \textbf{SIH26069} & \textbf{6.4/10} & Moderate & National Weather Big Data Analytics Platform... \\

21 & \textbf{SIH26191} & \textbf{5.9/10} & Moderate & Intelligent Identification of Hazard-Based Red Zones, Carrying Ca... \\

22 & \textbf{SIH26060} & \textbf{5.6/10} & Moderate & Digital Platform for efficient remote management of Indian Antarc... \\

23 & \textbf{SIH26074} & \textbf{5.6/10} & Moderate & Downscaling of weather forecast from Block level to Panchayat lev... \\

24 & \textbf{SIH26206} & \textbf{5.6/10} & Moderate & Student Innovation... \\

\bottomrule
\end{tabular}

\vspace{0.8cm}
\newpage

\tableofcontents
\newpage

\section{Detailed Problem Statements}

\needspace{6\baselineskip}
\phantomsection
\addcontentsline{toc}{subsection}{[9.2/10] [SIH26084] Convective scale nowcasting for Thunderstorms, Hail \& Cloudbursts (06 hr)}

\begin{tcolorbox}[pscard, title={"\textbf{SIH26084} \quad $\vert$ \quad Convective scale nowcasting for Thunderstorms, Hail \& Cloudbursts (06 hr)}]
\textbf{\large Convective scale nowcasting for Thunderstorms, Hail \& Cloudbursts (06 hr)}\\[0.25cm]

\begin{tabular}{@{}lp{12.5cm}@{}}
\textbf{Difficulty Rating:} & \textcolor{diffextreme}{\textbf{\Large 9.2 / 10}} \quad (\textbf{Extreme (Expert Level)}) \\
\textbf{Problem ID:} & 26084 \quad $\vert$ \quad \textbf{Category:} \textbf{Software} \quad $\vert$ \quad \textbf{Theme:} \textbf{Disaster Management} \\
\textbf{Organization:} & Ministry of Earth Sciences (MoES) \\
\textbf{Department:} & National Centre for Medium Range Weather Forecasting (NCMRWF) \\
\textbf{Recommended Stack:} & \texttt{PyTorch / TensorFlow, GDAL / Rasterio, Xarray, NetCDF4 / cfgrib, GeoPandas, PostGIS} \\
\textbf{Key Challenges:} & \textit{\small Physics-informed / Numerical simulation \& fluid dynamics coupling; High-volume Doppler radar \& satellite raster/GRIB2 ingestion; Ultra-low-latency nowcasting \& instant push alert pipelines; Specialized meteorological/hydrological physics domain modeling} \\
\end{tabular}

\vspace{0.2cm}
\hrule
\vspace{0.2cm}

\textbf{\color{sihnavy}Problem Statement Scope \& Requirements:}\\[0.15cm]
\noindent\small Convective storms, such as severe thunderstorms, hail, downburst winds, and cloudbursts are among India's deadliest natural hazards, especially during the pre-monsoon and monsoon seasons.Despite advancements in Numerical Weather Prediction (NWP) models, traditional systems often fail to accurately capture these mesoscale extreme weather events.The primary limitation stems from spatial and temporal constraints: these violent storms develop rapidly within a window of minutes and occur at localized scales that slip through coarse grid resolutions. Current early warning infrastructures struggle to provide high-resolution, short-term forecasts (0 - 6 hours), leaving local administrations, aviation sectors, and rural farming communities vulnerable to sudden, devastating impacts. The challenge is to build a real-time, convective-scale Nowcasting System (0 - 6 hour lead time) operating at a hyper-local 1 - 3 km spatial resolution. Because traditional physics-based models are too computationally slow to simulate these rapid developments in real-time, participants must design a system rooted in Multi-Source Data Fusion architectures. The core objective is to ingest high-frequency, heterogeneous meteorological streams, automatically detect early convective initiation, and dynamically forecast severe storm parameters (including lightning strike density, hail probability, downburst velocity, and cloudburst thresholds).Design a robust, real-time ingestion engine that fuses data streams from multiple sources: Doppler Weather Radars (DWR - reflectivity and velocity fields), geostationary satellite imagery (INSAT-3D/3DR thermal/infrared bands), and ground-based lightning detection networks. A real-time,interactive GIS-mapped dashboard showcasing high-resolution (1 - 3 km) hazard zones with live countdown clocks for storm arrivals.\par\vspace{0.06cm}

\end{tcolorbox}
\vspace{0.4cm}
\needspace{6\baselineskip}
\phantomsection
\addcontentsline{toc}{subsection}{[9.2/10] [SIH26085] Urban Flood Nowcasting System (Drainage and Rainfall Coupling)}

\begin{tcolorbox}[pscard, title={"\textbf{SIH26085} \quad $\vert$ \quad Urban Flood Nowcasting System (Drainage and Rainfall Coupling)}]
\textbf{\large Urban Flood Nowcasting System (Drainage and Rainfall Coupling)}\\[0.25cm]

\begin{tabular}{@{}lp{12.5cm}@{}}
\textbf{Difficulty Rating:} & \textcolor{diffextreme}{\textbf{\Large 9.2 / 10}} \quad (\textbf{Extreme (Expert Level)}) \\
\textbf{Problem ID:} & 26085 \quad $\vert$ \quad \textbf{Category:} \textbf{Software} \quad $\vert$ \quad \textbf{Theme:} \textbf{Disaster Management} \\
\textbf{Organization:} & Ministry of Earth Sciences (MoES) \\
\textbf{Department:} & National Centre for Medium Range Weather Forecasting (NCMRWF) \\
\textbf{Recommended Stack:} & \texttt{PyTorch / TensorFlow, GDAL / Rasterio, Xarray, NetCDF4 / cfgrib, GeoPandas, PostGIS} \\
\textbf{Key Challenges:} & \textit{\small Physics-informed / Numerical simulation \& fluid dynamics coupling; High-volume Doppler radar \& satellite raster/GRIB2 ingestion; Ultra-low-latency nowcasting \& instant push alert pipelines; Specialized meteorological/hydrological physics domain modeling} \\
\end{tabular}

\vspace{0.2cm}
\hrule
\vspace{0.2cm}

\textbf{\color{sihnavy}Problem Statement Scope \& Requirements:}\\[0.15cm]
\noindent\small Urban flooding in major Indian metros like Mumbai, Delhi, and Chennai has become an annual crisis. Traditional Numerical Weather Prediction (NWP) models fall short because knowing how much rain will fall does not automatically translate into knowing where the streets will flood.Urban flooding is a hyper-local phenomenon dictated by micro-topography, concrete imperviousness, and heavily strained, invisible drainage networks. Currently, municipal bodies lack real-time, street-level predictive systems. Consequently, cities are caught off guard by rapid water accumulation, leading to severe traffic gridlocks, economic disruption, and loss of life. The challenge is to design a high-resolution, real-time Urban Flood Nowcasting System (0 - 3 hour lead time) capable of predicting street-level inundation before it happens. Participants must move away from isolated weather models and instead build a coupled framework. This system must fuse real-time rainfall nowcasts with high-resolution Digital Elevation Models (DEM) and a graph-based mathematical model of the city's underground drainage network. By mapping how water flows, accumulates, and surcharges across concrete surfaces and drainage nodes, the solution should pinpoint exactly which streets or intersections will flood.Develop a pipeline that takes high-resolution rainfall nowcasts (from Doppler Weather Radars) and instantly routes that volume across a 2D surface terrain model. Represent the city's stormwater drain network as a directed graph (nodes as manholes/inlets, edges as pipes/canals). The model must calculate hydraulic capacity and predict where blockages or overcapacity will cause backflow onto the streets. A dynamic, web-based GIS dashboard showing real-time, street-by-street flooding projections (e.g., water depth estimations in centimeters) with a 0 - 3 hour forward-looking window.An API utility that can interface with navigation maps to suggest flood-safe alternative routes for emergency services, public transit, and commuters during heavy downpours.\par\vspace{0.06cm}

\end{tcolorbox}
\vspace{0.4cm}
\needspace{6\baselineskip}
\phantomsection
\addcontentsline{toc}{subsection}{[8.6/10] [SIH26077] AI-Driven Hyper-Local Early Warning System for Severe Weather Nowcasting}

\begin{tcolorbox}[pscard, title={"\textbf{SIH26077} \quad $\vert$ \quad AI-Driven Hyper-Local Early Warning System for Severe Weather Nowcasting}]
\textbf{\large AI-Driven Hyper-Local Early Warning System for Severe Weather Nowcasting}\\[0.25cm]

\begin{tabular}{@{}lp{12.5cm}@{}}
\textbf{Difficulty Rating:} & \textcolor{diffextreme}{\textbf{\Large 8.6 / 10}} \quad (\textbf{Extreme (Expert Level)}) \\
\textbf{Problem ID:} & 26077 \quad $\vert$ \quad \textbf{Category:} \textbf{Software} \quad $\vert$ \quad \textbf{Theme:} \textbf{Disaster Management} \\
\textbf{Organization:} & Ministry of Earth Sciences (MoES) \\
\textbf{Department:} & National Centre for Medium Range Weather Forecasting (NCMRWF) \\
\textbf{Recommended Stack:} & \texttt{PyTorch / TensorFlow, GDAL / Rasterio, Xarray, NetCDF4 / cfgrib, GeoPandas, PostGIS} \\
\textbf{Key Challenges:} & \textit{\small Physics-informed / Numerical simulation \& fluid dynamics coupling; High-volume Doppler radar \& satellite raster/GRIB2 ingestion; Ultra-low-latency nowcasting \& instant push alert pipelines; Disaster response planning \& resource management} \\
\end{tabular}

\vspace{0.2cm}
\hrule
\vspace{0.2cm}

\textbf{\color{sihnavy}Problem Statement Scope \& Requirements:}\\[0.15cm]
\noindent\hspace*{1.2em}\hangindent=2.4em\hangafter=1\textcolor{sihblue}{$\bullet$}~\small Problem Statement India is highly vulnerable to rapidly intensifying, localized extreme weather events such as cloudbursts, severe thunderstorms, and flash floods. Traditional physics-based Numerical Weather Prediction (NWP) models often suffer from computational latency and struggle to capture the rapid, small-scale atmospheric changes that preceded these events. There is a critical need for a real-time, hyper-local early warning system capable of 'nowcasting' severe weather 2 to 6 hours before impact, providing actionable lead time for disaster management.\par\vspace{0.05cm}
\noindent\hspace*{1.2em}\hangindent=2.4em\hangafter=1\textcolor{sihblue}{$\bullet$}~\small Proposed Solution We propose an advanced AI predictive engine designed for high-precision severe-weather nowcasting. Specifically, the system simultaneously predicts the onset of highly localized, rapidly intensifying events, namely severe thunderstorms, cloudbursts, and the subsequent flash floods,with an actionable lead time of 2 to 6 hours. Instead of relying on computationally intensive thermodynamic simulations, the system utilizes a spatiotemporal deep learning architecture to recognize the complex, multivariate atmospheric signatures that precede these extreme events. A critical component of this methodology is storm nowcasting using variations in integrated water vapor (IWV). By tracking rapid spatial and temporal accumulations of IWV, the model accurately identifies the concentrated moisture pools required for heavy precipitation. To predict multiple extreme events simultaneously, the engine employs a multi-task learning approach. A shared neural network backbone extracts foundational atmospheric features (moisture, instability, and lift) from the input grids. The network then branches into distinct output layers, allowing a single unified model to generate hyper-local probability risk maps for thunderstorms, cloudbursts, and flash floods simultaneously, entirely bypassing the computational latency typical of traditional numerical weather prediction (NWP) models. Predictive Matrix: Key Atmospheric Variables Severe convective storms require three primary ingredients: moisture, instability, and lift. Our AI model tracks the critical precursors across all three categories to ensure high accuracy and low false-alarm rates: ? Moisture Availability (The Fuel): The cornerstone of our storm nowcasting is the capture of integrated water vapor (IWV) variations. By tracking rapid spatial and temporal accumulations of IWV from satellites, the model identifies the concentrated moisture pools that trigger localized cloudbursts. ? Atmospheric Instability (The Energy): The model assesses the atmosphere's thermal profile to determine if it is buoyant enough to support explosive vertical cloud growth. High Convective Available Potential Energy (CAPE) paired with eroding Convective Inhibition (CIN) serves as a prime indicator of impending severe thunderstorms. Kinematics and Lift (The Trigger \& Structure): Low-level convergence (wind vectors colliding at the surface) forces air upward, initiating the development of a storm cell. Furthermore, tracking vertical wind shear (changes in wind speed/direction with altitude) helps the model predict whether a storm will move quickly or remain stationary. Observational Signatures: Rapid cooling of cloud tops, measured as the Cloud Top Temperature(CTT) Drop Rate, provides real-time validation of explosive vertical updrafts within the system. Topographic Dynamics (The Flood Catalyst): To accurately predict flash floods, the AI overlays the atmospheric probability maps onto a high-resolution Digital Elevation Model (DEM). This allows the system to calculate how terrain slope, elevation, and natural drainage basins will channel the extreme precipitation generated by a predicted cloudburst.To capture these predictors with hyper-local accuracy, the model fuses multi-modal, high resolution datasets: IMDAA Reanalysis Data (Historical Baseline \& Thermodynamics): Multi-level air temperature, specific humidity profiles (for calculating CAPE/CIN), geopotential height, and U/V wind components (for calculating shear and convergence). Satellite Observations (INSAT-3D/3DR via MOSDAC): Water Vapor (WV) Channels. This is essential for deriving real-time Integrated Water Vapor (IWV) fluctuations necessary for our storm nowcasting. Thermal Infrared (TIR) Channels: Utilized to calculate the rapid Cloud Top Temperature (CTT) drop rate. Quantitative Precipitation Estimation (QPE): Satellite-derived precipitation estimates are used to monitor real-time rainfall intensity, serving as a reliable, openly accessible alternative to ground-based radar. Digital Elevation Model (DEM): High-resolution topographical data (such as ISRO's CartoDEM or SRTM) provides a static baseline of elevation, slope, and surface drainage networks, enabling translation of atmospheric cloudburst predictions into actionable flash flood warnings on the ground.\par\vspace{0.05cm}
\noindent\hspace*{1.2em}\hangindent=2.4em\hangafter=1\textcolor{sihblue}{$\bullet$}~\small Technical Methodology ? Data Fusion \& Alignment: Raw data from IMDAA reanalysis, INSAT-3D/3DR satellite observations, and high-resolution Digital Elevation Models (DEM) are ingested, normalized, and mapped onto a unified spatiotemporal grid (e.g., using multi-dimensional array structures). This ensures that all dynamic atmospheric predictors - such as specific humidity and cloud-top temperatures - and static surface variables align geographically and chronologically for seamless multimodal processing. ? Multi-Variate Feature Extraction \& Multi-Task Inference: A shared multi-modal spatiotemporal transformer network continuously analyzes real-time satellite grids, specifically tracking critical IWV variations and CTT drop rates, against the IMDAA-derived thermodynamic baselines using cross-attention mechanisms. Utilizing a Multi-Task Learning (MTL) architecture,the network branches into distinct output 'heads.' This allows the unified model to simultaneously process the aligned data and generate distinct, hyper-local probability maps for severe thunderstorms, cloudbursts, and flash floods without computational bottlenecking. ? Automated Alerting: When the predictive matrix breaches the signature thresholds of a severe event, the engine generates a spatial risk map and pushes automated, categorized alerts via a lightweight API.\par\vspace{0.05cm}
\noindent\hspace*{1.2em}\hangindent=2.4em\hangafter=1\textcolor{sihblue}{$\bullet$}~\small Expected Solution The final deliverable for the Smart India Hackathon will be a fully functional, real-time prototypeof the AI-Driven Hyper-Local Early Warning System. At its core is a deployed multi-task inference engine that continuously ingests live INSAT satellite data and IMDAA thermodynamic baselines to simultaneously generate predictive risk maps for severe thunderstorms, cloudbursts, and flash floods within a 2 to 6-hour predictive window. This backend integrates with an interactive, webbased spatial dashboard designed for disaster management authorities, featuring dynamic risk maps overlaid on a Digital Elevation Model (DEM) and an Explainable AI (XAI) module that transparently displays meteorological triggers. Finally, an automated API will translate these predictive insights into immediate, categorized alerts sent directly to first responders and vulnerable communities the moment critical thresholds are breached.\par\vspace{0.05cm}

\end{tcolorbox}
\vspace{0.4cm}
\needspace{6\baselineskip}
\phantomsection
\addcontentsline{toc}{subsection}{[8.6/10] [SIH26161] Dam Break Inundation Modelling Using Hydrodynamic Modelling of any River}

\begin{tcolorbox}[pscard, title={"\textbf{SIH26161} \quad $\vert$ \quad Dam Break Inundation Modelling Using Hydrodynamic Modelling of any River}]
\textbf{\large Dam Break Inundation Modelling Using Hydrodynamic Modelling of any River}\\[0.25cm]

\begin{tabular}{@{}lp{12.5cm}@{}}
\textbf{Difficulty Rating:} & \textcolor{diffextreme}{\textbf{\Large 8.6 / 10}} \quad (\textbf{Extreme (Expert Level)}) \\
\textbf{Problem ID:} & 26161 \quad $\vert$ \quad \textbf{Category:} \textbf{Software} \quad $\vert$ \quad \textbf{Theme:} \textbf{Disaster Management} \\
\textbf{Organization:} & National Technical Research Organisation (NTRO) \\
\textbf{Department:} & National Technical Research Organisation (NTRO) \\
\textbf{Recommended Stack:} & \texttt{PyTorch / TensorFlow, GDAL / Rasterio, Xarray, NetCDF4 / cfgrib, GeoPandas, PostGIS} \\
\textbf{Key Challenges:} & \textit{\small Physics-informed / Numerical simulation \& fluid dynamics coupling; High-volume Doppler radar \& satellite raster/GRIB2 ingestion; Cloud web dashboard \& RESTful API microservices; Specialized meteorological/hydrological physics domain modeling} \\
\end{tabular}

\vspace{0.2cm}
\hrule
\vspace{0.2cm}

\textbf{\color{sihnavy}Problem Statement Scope \& Requirements:}\\[0.15cm]
\noindent\hspace*{1.2em}\hangindent=2.4em\hangafter=1\textcolor{sihblue}{$\bullet$}~\small Background In India, due to natural disaster various natural dam / lake formations were observed which can be a major reason of flash flood in the lower catchment, for example, natural lake formed over the Rishi Ganga river of Uttarakhand in Feb 2021, Wapriyang river in Nov 2021, Phuktal river near Sumdo, J\&K in Mar 15, Kosi river in 2008 etc. Devastating flood happened in the Kashmir valley, Assam in 2014 and many other places over a period of time. Therefore, simulation modelling for flash flood and scenario generation is important from Humanitarian Assistance and Disaster Relief (HADR) point of view. Another important aspect is water release issues from the Dam of major rivers. In the crisis situation, if the dam brakes, how much water will flow into the river and what are the area it will inundated / impacted need to be estimated. In order to carry out this work simulation modelling needs to be done for the same. In other way if any dam break situation happened then what will be the impacted area. Development of a modelling framework is required which will carry out simulation modelling for the same.\par\vspace{0.05cm}
\noindent\hspace*{1.2em}\hangindent=2.4em\hangafter=1\textcolor{sihblue}{$\bullet$}~\small Description The above problem statement envisages that a software tool need to be developed which should automatically carry out the simulation modelling for Dam break analysis and identify the inundated area due to flash flood in the lower catchment. The modelling framework should be developed using hydrological data, DEM and satellite imagery of any river. The software/ tools should be capable of carrying out the simulation modelling of water flow in case of dam break or water release through 'Smooth Particle Hydrodynamics' and 'Delf3D' model and compare the scenario.\par\vspace{0.05cm}
\noindent\hspace*{1.2em}\hangindent=2.4em\hangafter=1\textcolor{sihblue}{$\bullet$}~\small Expected Solution/Deliverables The proposed study aims to illustrate the current problems regarding framework generation of Humanitarian Assistance and Disaster Relief using simulation modelling related to flood management as follows:\par\vspace{0.05cm}
\noindent\hangindent=1.8em\hangafter=1\textbf{i.}~\small Creation of generalized modelling framework to predict / simulate dam break/ river blockage analysis providing the necessary inputs on the basis of sudden water surge as well as loss and damage analysis using 'Smooth Particle Hydrodynamics model and Delf3D model'. ii. Building a customized tool/ framework so that it is possible to generate a flood inundation simulation scenario using different input datasets. iii. Developing a Dashboard for providing modelling input and output visualization framework (GUI). The program should support the large volume of data. Output should be converted to .shp or .Kml file. iv. Additionally, developing a framework for near real time flood analysis through Google Earth Engine with the help of open source data.\par\vspace{0.06cm}
\noindent\hangindent=1.8em\hangafter=1\textbf{v.}~\small Simulation needs to be done by taking the any river and Dam data (open source) of India during the final demonstration of the software.\par\vspace{0.06cm}


\vspace{0.2cm}
\hrule
\vspace{0.2cm}
{\footnotesize
\textbf{Dataset / Resources:} Open source Remote Sensing data (Sentinel, Landsat or any open source satellite image) and ASTER/ STRM or any other DEM.
}

\end{tcolorbox}
\vspace{0.4cm}
\needspace{6\baselineskip}
\phantomsection
\addcontentsline{toc}{subsection}{[8.3/10] [SIH26072] AIML based Nowcasting of thunderstorm and lightning using atmospheric obser...}

\begin{tcolorbox}[pscard, title={"\textbf{SIH26072} \quad $\vert$ \quad AIML based Nowcasting of thunderstorm and lightning using atmospheric obser...}]
\textbf{\large AIML based Nowcasting of thunderstorm and lightning using atmospheric observation including multiple radars, satellite, lightning and model data.}\\[0.25cm]

\begin{tabular}{@{}lp{12.5cm}@{}}
\textbf{Difficulty Rating:} & \textcolor{diffveryhard}{\textbf{\Large 8.3 / 10}} \quad (\textbf{Very Hard (Advanced)}) \\
\textbf{Problem ID:} & 26072 \quad $\vert$ \quad \textbf{Category:} \textbf{Software} \quad $\vert$ \quad \textbf{Theme:} \textbf{Disaster Management} \\
\textbf{Organization:} & Ministry of Earth Sciences (MoES) \\
\textbf{Department:} & India Meteorological Department \\
\textbf{Recommended Stack:} & \texttt{PyTorch / TensorFlow, GDAL / Rasterio, Xarray, NetCDF4 / cfgrib, FastAPI, Celery / Redis} \\
\textbf{Key Challenges:} & \textit{\small Advanced Spatio-temporal AI / Multi-modal Deep Learning models; High-volume Doppler radar \& satellite raster/GRIB2 ingestion; Ultra-low-latency nowcasting \& instant push alert pipelines; Disaster response planning \& resource management} \\
\end{tabular}

\vspace{0.2cm}
\hrule
\vspace{0.2cm}

\textbf{\color{sihnavy}Problem Statement Scope \& Requirements:}\\[0.15cm]
\noindent\small AIML based Nowcasting of thunderstorm and lightning using atmospheric observation including multiple radars, satellite, lightning and model data.\par\vspace{0.06cm}

\end{tcolorbox}
\vspace{0.4cm}
\needspace{6\baselineskip}
\phantomsection
\addcontentsline{toc}{subsection}{[8.2/10] [SIH26001] AI-Based early warning and landslide Risk Monitoring System in NER}

\begin{tcolorbox}[pscard, title={"\textbf{SIH26001} \quad $\vert$ \quad AI-Based early warning and landslide Risk Monitoring System in NER}]
\textbf{\large AI-Based early warning and landslide Risk Monitoring System in NER}\\[0.25cm]

\begin{tabular}{@{}lp{12.5cm}@{}}
\textbf{Difficulty Rating:} & \textcolor{diffveryhard}{\textbf{\Large 8.2 / 10}} \quad (\textbf{Very Hard (Advanced)}) \\
\textbf{Problem ID:} & 26001 \quad $\vert$ \quad \textbf{Category:} \textbf{Software} \quad $\vert$ \quad \textbf{Theme:} \textbf{Disaster Management} \\
\textbf{Organization:} & Ministry of Development of North Eastern Region (MDoNER) \\
\textbf{Department:} & Ministry of Development of North Eastern Region (MDoNER) \\
\textbf{Recommended Stack:} & \texttt{PyTorch / TensorFlow, GDAL / Rasterio, Xarray, NetCDF4 / cfgrib, GeoPandas, PostGIS} \\
\textbf{Key Challenges:} & \textit{\small Predictive ML \& anomaly detection models; High-volume Doppler radar \& satellite raster/GRIB2 ingestion; Ultra-low-latency nowcasting \& instant push alert pipelines; Geotechnical vulnerability \& multi-hazard spatial zoning} \\
\end{tabular}

\vspace{0.2cm}
\hrule
\vspace{0.2cm}

\textbf{\color{sihnavy}Problem Statement Scope \& Requirements:}\\[0.15cm]
\vspace{0.18cm}\noindent\textbf{\color{sihnavy}Background:}~\small The North Eastern Region (NER) frequently faces landslides, flash floods, road blockages, and slope failures due to heavy rainfall, fragile terrain, and unplanned hill cutting. These incidents often disrupt connectivity, damage infrastructure, delay emergency response, and isolate remote villages for days. Currently, monitoring of vulnerable zones is mostly reactive and dependent on manual reporting. There is limited use of real-time predictive systems for identifying high-risk zones and issuing early warnings to authorities and local communities. With increasing climate vulnerability in the region, there is a need for an AI-enabled real-time monitoring and prediction system that can help authorities take preventive action before disasters occur.\par\vspace{0.06cm}
\vspace{0.18cm}\noindent\textbf{\color{sihnavy}Description:}\par\vspace{0.06cm}
\noindent\small This problem statement proposes the development of an Al-powered early warning and monitoring platform capable of predicting and tracking landslide-prone areas in real time across the North Eastern Region.\par\vspace{0.06cm}
\vspace{0.18cm}\noindent\textbf{\color{sihnavy}The solution should:}\par\vspace{0.06cm}
\noindent\hangindent=1.8em\hangafter=1\textbf{a.}~\small Collect and analyse data from: Rainfall patterns Soil moisture sensors Satellite imagery Terrain/slope data Historical landslide records\par\vspace{0.06cm}
\noindent\hangindent=1.8em\hangafter=1\textbf{b.}~\small Use AI/ML models to identify high-risk zones and predict possible landslide events.\par\vspace{0.06cm}
\noindent\hangindent=1.8em\hangafter=1\textbf{c.}~\small Provide real-time alerts to district administrations, disaster management authorities, and local communities.\par\vspace{0.06cm}
\noindent\hangindent=1.8em\hangafter=1\textbf{d.}~\small Integrate GIS mapping for visualization of vulnerable roads, villages, and infrastructure.\par\vspace{0.06cm}
\noindent\hangindent=1.8em\hangafter=1\textbf{e.}~\small Allow citizens/field officials to upload geo-tagged photos/videos of cracks, slope movement or blocked roads.\par\vspace{0.06cm}
\noindent\hangindent=1.8em\hangafter=1\textbf{f.}~\small Generate dashboards showing:\par\vspace{0.06cm}
\noindent\hspace*{1.2em}\hangindent=2.4em\hangafter=1\textcolor{sihblue}{$\bullet$}~\small Risk severity levels\par\vspace{0.05cm}
\noindent\hspace*{1.2em}\hangindent=2.4em\hangafter=1\textcolor{sihblue}{$\bullet$}~\small Road connectivity status\par\vspace{0.05cm}
\noindent\hspace*{1.2em}\hangindent=2.4em\hangafter=1\textcolor{sihblue}{$\bullet$}~\small Weather-linked risk forecasts\par\vspace{0.05cm}
\noindent\hspace*{1.2em}\hangindent=2.4em\hangafter=1\textcolor{sihblue}{$\bullet$}~\small Emergency response prioritisation. Support multilingual notifications and low-network/offline functionality for remote areas.\par\vspace{0.05cm}
\vspace{0.18cm}\noindent\textbf{\color{sihnavy}Expected Solution:}\par\vspace{0.06cm}
\noindent\small A scalable Al-based software platform with:\par\vspace{0.06cm}
\noindent\hspace*{1.2em}\hangindent=2.4em\hangafter=1\textcolor{sihblue}{$\bullet$}~\small Real-time GIS dashboard and risk heatmaps\par\vspace{0.05cm}
\noindent\hspace*{1.2em}\hangindent=2.4em\hangafter=1\textcolor{sihblue}{$\bullet$}~\small AI/ML-based predictive analytics engine\par\vspace{0.05cm}
\noindent\hspace*{1.2em}\hangindent=2.4em\hangafter=1\textcolor{sihblue}{$\bullet$}~\small Mobile/web application for field reporting and alerts.\par\vspace{0.05cm}
\noindent\hspace*{1.2em}\hangindent=2.4em\hangafter=1\textcolor{sihblue}{$\bullet$}~\small Integration with IMD weather APIs, satellite feeds, and sensor data\par\vspace{0.05cm}
\noindent\hspace*{1.2em}\hangindent=2.4em\hangafter=1\textcolor{sihblue}{$\bullet$}~\small Automated SMS/app-based early warning system\par\vspace{0.05cm}
\noindent\hspace*{1.2em}\hangindent=2.4em\hangafter=1\textcolor{sihblue}{$\bullet$}~\small Cloud-based architecture with offline sync support for remote regions The solution should improve disaster preparedness, reduce loss of life and infrastructure damage, and strengthen climate-resilient governance in the North Eastern Region.\par\vspace{0.05cm}

\end{tcolorbox}
\vspace{0.4cm}
\needspace{6\baselineskip}
\phantomsection
\addcontentsline{toc}{subsection}{[8.0/10] [SIH26071] AI/ML-Based Integrated heavy rainfall Early Warning and Inundation Predicti...}

\begin{tcolorbox}[pscard, title={"\textbf{SIH26071} \quad $\vert$ \quad AI/ML-Based Integrated heavy rainfall Early Warning and Inundation Predicti...}]
\textbf{\large AI/ML-Based Integrated heavy rainfall Early Warning and Inundation Prediction System using Satellite, Radar, observational Weather and numerical weather prediction model data.}\\[0.25cm]

\begin{tabular}{@{}lp{12.5cm}@{}}
\textbf{Difficulty Rating:} & \textcolor{diffveryhard}{\textbf{\Large 8.0 / 10}} \quad (\textbf{Very Hard (Advanced)}) \\
\textbf{Problem ID:} & 26071 \quad $\vert$ \quad \textbf{Category:} \textbf{Software} \quad $\vert$ \quad \textbf{Theme:} \textbf{Disaster Management} \\
\textbf{Organization:} & Ministry of Earth Sciences (MoES) \\
\textbf{Department:} & India Meteorological Department \\
\textbf{Recommended Stack:} & \texttt{PyTorch / TensorFlow, GDAL / Rasterio, Xarray, NetCDF4 / cfgrib, FastAPI / Django, PostgreSQL} \\
\textbf{Key Challenges:} & \textit{\small Physics-informed / Numerical simulation \& fluid dynamics coupling; High-volume Doppler radar \& satellite raster/GRIB2 ingestion; Cloud web dashboard \& RESTful API microservices; Disaster response planning \& resource management} \\
\end{tabular}

\vspace{0.2cm}
\hrule
\vspace{0.2cm}

\textbf{\color{sihnavy}Problem Statement Scope \& Requirements:}\\[0.15cm]
\noindent\small AI/ML-Based Integrated heavy rainfall Early Warning and Inundation Prediction System using Satellite, Radar, observational Weather and numerical weather prediction model data.\par\vspace{0.06cm}

\end{tcolorbox}
\vspace{0.4cm}
\needspace{6\baselineskip}
\phantomsection
\addcontentsline{toc}{subsection}{[8.0/10] [SIH26078] AI-Driven Spatio-Temporal Tracking of Extreme Weather Anomalies in Medium-R...}

\begin{tcolorbox}[pscard, title={"\textbf{SIH26078} \quad $\vert$ \quad AI-Driven Spatio-Temporal Tracking of Extreme Weather Anomalies in Medium-R...}]
\textbf{\large AI-Driven Spatio-Temporal Tracking of Extreme Weather Anomalies in Medium-Range Forecasts}\\[0.25cm]

\begin{tabular}{@{}lp{12.5cm}@{}}
\textbf{Difficulty Rating:} & \textcolor{diffveryhard}{\textbf{\Large 8.0 / 10}} \quad (\textbf{Very Hard (Advanced)}) \\
\textbf{Problem ID:} & 26078 \quad $\vert$ \quad \textbf{Category:} \textbf{Software} \quad $\vert$ \quad \textbf{Theme:} \textbf{Disaster Management} \\
\textbf{Organization:} & Ministry of Earth Sciences (MoES) \\
\textbf{Department:} & National Centre for Medium Range Weather Forecasting (NCMRWF) \\
\textbf{Recommended Stack:} & \texttt{PyTorch, Scikit-Learn, XGBoost, Pandas, GeoPandas, PostGIS} \\
\textbf{Key Challenges:} & \textit{\small Physics-informed / Numerical simulation \& fluid dynamics coupling; High-volume Doppler radar \& satellite raster/GRIB2 ingestion; Cloud web dashboard \& RESTful API microservices; Disaster response planning \& resource management} \\
\end{tabular}

\vspace{0.2cm}
\hrule
\vspace{0.2cm}

\textbf{\color{sihnavy}Problem Statement Scope \& Requirements:}\\[0.15cm]
\noindent\hspace*{1.2em}\hangindent=2.4em\hangafter=1\textcolor{sihblue}{$\bullet$}~\small Problem Statement Identifying and tracking the exact geographic footprints of extreme weather anomalies (such as severe cyclones, heat domes, or cold waves) within massive global Numerical Weather Prediction (NWP) outputs is computationally intensive and heavily reliant on manual interpretation. In medium-range forecasting (3 to 10 days), atmospheric chaos renders traditional deterministic models highly uncertain. Furthermore, standard deep learning models (like standard CNNs or U-Nets) suffer from spectral smoothing - they tend to 'average out' spatial data, which destroys the extreme amplitudes (the high-intensity peaks of rainfall or wind speed) that forecasters actually need to track. There is a critical gap between broad, coarse 12 km global ensemble datasets and localized, high-fidelity threat tracking.\par\vspace{0.05cm}
\noindent\hspace*{1.2em}\hangindent=2.4em\hangafter=1\textcolor{sihblue}{$\bullet$}~\small Proposed Solution We propose an automated, state-of-the-art AI tracking and downscaling pipeline that shifts the paradigm from manual weather data sorting to automated, physics-informed anomaly tracking.Instead of relying on a single deterministic forecast run, our system directly processes multivariable, 4D Ensemble Prediction Systems (EPS) data.The system uses a two-stage hybrid AI architecture to solve the spectral smoothing problem:First, it utilizes a graph neural network (GNN) to map atmospheric variables onto a spherical mesh,instantly isolating moving anomalies and calculating their trajectory over a 3- to 10-day forecast window. Second, it pipes this isolated region into a generative diffusion model to perform statistical downscaling. This physics-constrained generative model mathematically derives a hyper-local 5km subgrid impact zone without flattening or blurring the severe amplitudes of the extreme weather event.\par\vspace{0.05cm}
\noindent\hspace*{1.2em}\hangindent=2.4em\hangafter=1\textcolor{sihblue}{$\bullet$}~\small Technical Methodology \& Architecture Spherical Anomaly Tracking (Stage 1 GNN): To eliminate the geographic distortions caused by processing the spherical Earth on flat 2D pixel grids, the system maps the 12 km NCMRWF Global Ensemble (NEPS-G) grids directly onto an icosahedral mesh. The message-passing GNN calculates the Extreme Forecast Index (EFI) against a 30-year historical ERA5 baseline distribution to isolate standard deviations and draw a macro-scale temporal bounding box around the anomaly's trajectory. Amplitude-Preserving Downscaling (Stage 2 Diffusion): The system passes the cropped, macroscale bounding box into a conditional denoising diffusion probabilistic model. Rather than optimizing for mean errors (which blurs peaks), the diffusion model learns the physical relationships between synoptic-scale features and regional topography. It iteratively generates high-resolution, high-amplitude local weather scenarios, downscaling the 12 km grid into a 5 km grid. Physics-Informed Constraints: To ensure the model remains scientifically accurate, we embed fluid dynamics and thermodynamic conservation laws directly into the neural network's loss function. The model is mathematically penalized if it generates physically impossible weather states (e.g., severe downpours missing corresponding moisture convergence vectors).\par\vspace{0.05cm}
\noindent\hspace*{1.2em}\hangindent=2.4em\hangafter=1\textcolor{sihblue}{$\bullet$}~\small Datasets and Tools ? AI Frameworks: PyTorch / JAX (engineered with custom, physics-guided loss functions),Deep Graph Library (DGL) for icosahedral mesh networks, and Hugging Face Diffusers for generative downscaling. ? Data Wrangling \& Geospatial Tools: Xarray and Dask for processing parallelized, multigigabyte 4D NetCDF/GRIB2 arrays; MetPy for physical meteorological equations; Cartopy for geographical map projections. ? Training \& Testing Datasets: ? Baseline: Historical IMDAA / ERA5 reanalysis data to establish the climatological norm. ? Forecast Inputs: Historical NCUM (12 km deterministic) and NEPS-G (12 km global ensemble) datasets containing documented extreme historical events (e.g., Cyclone Amphan, severe North India heatwaves).\par\vspace{0.05cm}
\noindent\hspace*{1.2em}\hangindent=2.4em\hangafter=1\textcolor{sihblue}{$\bullet$}~\small Expected Outcome \& Key Deliverables The Tracking Core: A production-ready Spatio-Temporal GNN module that continuously processes global NWP streams to output dynamic, automated 4D bounding boxes around evolving weather threats. The Downscaling Core: A generative diffusion module capable of ingesting a 12 km resolution anomaly slice and outputting a probabilistically sound, 5 km resolution sub-grid array that retains extreme value amplitudes. The Visualization \& Alert Dashboard: An automated system that translates the mathematical 5 km centroid arrays into clean, geographic visual layers. The Alerting API: A lightweight, production-ready REST API that programmatically drops a pinpoint coordinate at the core of the severe anomaly and triggers categorized spatial alerts (low, moderate, and severe) across a precise 5 km geographical impact radius.\par\vspace{0.05cm}
\noindent\hspace*{1.2em}\hangindent=2.4em\hangafter=1\textcolor{sihblue}{$\bullet$}~\small Use Cases \& Societal Impact Eliminating Alert Fatigue for the NDRF: Current weather alerts are often too broad, covering entire states or districts, which leads to public complacency. This solution allows meteorologists to issue hyper-localized, highly targeted warnings. It changes a generic'heavy rain in the district' alert into a precise 'high risk of flash flooding within your specific 5 km radius in the next 12 hours' alert, empowering first responders to deploy assets perfectly. Protecting Rural Economies: Grants farming communities a highly accurate, 3- to 10- day lead time regarding localized catastrophic anomalies like sudden frost, hail, or heat domes. This structural foresight lets farmers alter harvesting schedules or apply cropprotection covers, shielding rural livelihoods from sudden climate shocks. Democratizing Supercomputing Power: Once this hybrid AI pipeline is trained, it processes live inference data on a standard cloud GPU node in seconds, making high-fidelity climate forecasting highly affordable and easily accessible.\par\vspace{0.05cm}

\end{tcolorbox}
\vspace{0.4cm}
\needspace{6\baselineskip}
\phantomsection
\addcontentsline{toc}{subsection}{[7.9/10] [SIH26083] Extreme Heatwave Early Warning and Human Thermal Stress Index}

\begin{tcolorbox}[pscard, title={"\textbf{SIH26083} \quad $\vert$ \quad Extreme Heatwave Early Warning and Human Thermal Stress Index}]
\textbf{\large Extreme Heatwave Early Warning and Human Thermal Stress Index}\\[0.25cm]

\begin{tabular}{@{}lp{12.5cm}@{}}
\textbf{Difficulty Rating:} & \textcolor{diffveryhard}{\textbf{\Large 7.9 / 10}} \quad (\textbf{Very Hard (Advanced)}) \\
\textbf{Problem ID:} & 26083 \quad $\vert$ \quad \textbf{Category:} \textbf{Software} \quad $\vert$ \quad \textbf{Theme:} \textbf{Disaster Management} \\
\textbf{Organization:} & Ministry of Earth Sciences (MoES) \\
\textbf{Department:} & National Centre for Medium Range Weather Forecasting (NCMRWF) \\
\textbf{Recommended Stack:} & \texttt{PyTorch, Scikit-Learn, XGBoost, Pandas, GeoPandas, PostGIS} \\
\textbf{Key Challenges:} & \textit{\small Predictive ML \& anomaly detection models; High-volume Doppler radar \& satellite raster/GRIB2 ingestion; Cloud web dashboard \& RESTful API microservices; Specialized meteorological/hydrological physics domain modeling} \\
\end{tabular}

\vspace{0.2cm}
\hrule
\vspace{0.2cm}

\textbf{\color{sihnavy}Problem Statement Scope \& Requirements:}\\[0.15cm]
\noindent\small In recent years, the frequency, duration, and intensity of heatwaves across India have escalated sharply due to climate change. However, traditional meteorological warnings rely almost exclusively on ambient dry-bulb temperature thresholds. This creates a critical vulnerability:standard forecasts ignore the deadly compounding effects of relative humidity, wind speed, and solar radiation on the human body. A temperature of 40Â$^\circ$C at 20\% humidity feels vastly different from 40Â$^\circ$C at 70\% humidity - the latter can be fatal. Current public health infrastructure lacks localized, impact-based forecasting that translates raw weather data into actual physiological risk, human thermal stress levels, and projected mortality rates. The challenge is to design an intelligent, localized early warning system that shifts heatwave forecasting from 'what the weather will be' to 'what the weather will do' to human health. Participants need to build a predictive platform that computes a comprehensive Human Thermal Stress Index (integrating temperature, humidity, wind, and radiation) and links it directly to an automated Mortality Risk Index. The system should offer high-resolution forecasts to help municipal corporations, healthcare systems, and disaster management authorities deploy targeted,preemptive interventions. Develop algorithms to calculate advanced heat stress metrics such as the Wet-Bulb Globe Temperature (WBGT), Universal Thermal Climate Index (UTCI), or Heat Index (HI) rather than relying on temperature alone. Integrate historical public health, demographic (e.g., elderly or outdoor worker density), and localized weather data to predict heat-induced mortality and hospitalization spikes 3 to 5 days in advance. A dynamic GIS-mapped dashboard providing colorcoded, hyper-local alerts (Zone/Ward level) paired with actionable, automated public health advisories. An API capable of pushing automated SMS/WhatsApp regional alerts or localized triggers for city administration to initiate heat action plans (e.g., opening cooling centers, adjusting power grids, shifting outdoor work hours).\par\vspace{0.06cm}

\end{tcolorbox}
\vspace{0.4cm}
\needspace{6\baselineskip}
\phantomsection
\addcontentsline{toc}{subsection}{[7.3/10] [SIH26013] Automated lntegration and lntelligent Harmonization of Multi-source Geospat...}

\begin{tcolorbox}[pscard, title={"\textbf{SIH26013} \quad $\vert$ \quad Automated lntegration and lntelligent Harmonization of Multi-source Geospat...}]
\textbf{\large Automated lntegration and lntelligent Harmonization of Multi-source Geospatial Data for urban Land Record Management.}\\[0.25cm]

\begin{tabular}{@{}lp{12.5cm}@{}}
\textbf{Difficulty Rating:} & \textcolor{diffhard}{\textbf{\Large 7.3 / 10}} \quad (\textbf{Hard (Intermediate-High)}) \\
\textbf{Problem ID:} & 26013 \quad $\vert$ \quad \textbf{Category:} \textbf{Software} \quad $\vert$ \quad \textbf{Theme:} \textbf{Disaster Management} \\
\textbf{Organization:} & Ministry of Rural Development \\
\textbf{Department:} & Dept of land resources (DoLR) \\
\textbf{Recommended Stack:} & \texttt{PyTorch, Scikit-Learn, XGBoost, Pandas, GeoPandas, PostGIS} \\
\textbf{Key Challenges:} & \textit{\small Advanced Spatio-temporal AI / Multi-modal Deep Learning models; Multi-source geospatial GIS \& real-time sensor fusion; Cloud web dashboard \& RESTful API microservices; Disaster response planning \& resource management} \\
\end{tabular}

\vspace{0.2cm}
\hrule
\vspace{0.2cm}

\textbf{\color{sihnavy}Problem Statement Scope \& Requirements:}\\[0.15cm]
\vspace{0.18cm}\noindent\textbf{\color{sihnavy}Background:}~\small Urban land administration and cadastral management involve integration of multiple spatial and non-spatial datasets generated from various departments,agencies,and survey mechanisms. Under modern land governance programmes such as the NAKSHA Programme, large volumes of geospatial data are being generated through drone surveys, Orthorectified lmagery (ORl), DSM/DTM datasets, Ground Truthing (GT),GNSS surveys,municipal records,utility databases, and revenue land records. At present,harmonization and integration of these datasets largely depend on manual GIS workflows,which are time-consuming and prone to errors.With increasing availability of Al, GeoAl, and automated spatial processing technologies, there is significant scope for development of an intelligent system capable of automatically integrating and synchronizing multi-source geospatial datasets with feature-extracted cadastral data.\par\vspace{0.06cm}
\vspace{0.18cm}\noindent\textbf{\color{sihnavy}Description:}\par\vspace{0.06cm}
\noindent\small The proposed solution should develop an Al-enabled geospatial integration platform capable of automatically integrating, harmonizing, validating, and synchronizing multiple land-related datasets with Al-generated feature extraction outputs. The system should support integration of: . Drone imagery . Orthorectified lmagery (ORl) . DSM/DTM datasets . Existing cadastral maps . Revenue records . Municipal GIS layers . Utility network data . Ground Truthing (GT) datasets . GNSS/CORS survey data . Building footPrint datasets The solution should incorPorate: . Al/ML-based spatial matching algorithms . Automated topology correction . lntelligent attribute mapping . Geo-referencing and coordinate transformation engine . Change detection mechanisms . Spatial conflict resolution framework . Confidence scoring for integrated outputs\par\vspace{0.06cm}
\vspace{0.18cm}\noindent\textbf{\color{sihnavy}Expected Solution:}\par\vspace{0.06cm}
\noindent\small . The expected outcome is development of an intelligent geospatial integration framework capable of automatically harmonizing multi-source land-related datasets with Al-generated feature extraction outputs. .The final solution should: . Reduce manual GIS integration efforts . lmprove accuracy and consistency of urban land records . Enable seamless inter-departmental spatial data exchange . Accelerate cadastral finalization processes . lmprove interoperability of urban land information systems . Support standardized digital land governance Suggested Technologies: . Artificial lntelligence (Al) . Machine Learning (ML) . GeoAl . GIS \& Web-GlS . Spatial Databases . ETL Automation . Computer Vision . Cloud Computing . Spatial Analytics . API lntegration Frameworks\par\vspace{0.06cm}

\end{tcolorbox}
\vspace{0.4cm}
\needspace{6\baselineskip}
\phantomsection
\addcontentsline{toc}{subsection}{[7.3/10] [SIH26028] Dynamic Forecast of Expected Time of Arrival (ETA) for Coaching Trains}

\begin{tcolorbox}[pscard, title={"\textbf{SIH26028} \quad $\vert$ \quad Dynamic Forecast of Expected Time of Arrival (ETA) for Coaching Trains}]
\textbf{\large Dynamic Forecast of Expected Time of Arrival (ETA) for Coaching Trains}\\[0.25cm]

\begin{tabular}{@{}lp{12.5cm}@{}}
\textbf{Difficulty Rating:} & \textcolor{diffhard}{\textbf{\Large 7.3 / 10}} \quad (\textbf{Hard (Intermediate-High)}) \\
\textbf{Problem ID:} & 26028 \quad $\vert$ \quad \textbf{Category:} \textbf{Software} \quad $\vert$ \quad \textbf{Theme:} \textbf{Disaster Management} \\
\textbf{Organization:} & Ministry of Railways \\
\textbf{Department:} & Ministry of Railways \\
\textbf{Recommended Stack:} & \texttt{PyTorch, Scikit-Learn, XGBoost, Pandas, GeoPandas, PostGIS} \\
\textbf{Key Challenges:} & \textit{\small Predictive ML \& anomaly detection models; High-volume Doppler radar \& satellite raster/GRIB2 ingestion; Cloud web dashboard \& RESTful API microservices; Disaster response planning \& resource management} \\
\end{tabular}

\vspace{0.2cm}
\hrule
\vspace{0.2cm}

\textbf{\color{sihnavy}Problem Statement Scope \& Requirements:}\\[0.15cm]
\vspace{0.18cm}\noindent\textbf{\color{sihnavy}Background:}~\small Accurate forecasting of the Expected Time of Arrival (ETA) for coaching trains is vital for improving passenger satisfaction and operational efficiency in Indian Railways. Currently, ETA is often estimated using static schedules, current delays and in-built recovery times, which may not reflect real-time ground realities such as speed restrictions, congestion, unscheduled stoppages or historical patterns. As a result, passengers, station staff, and downstream logistics services face uncertainty and planning difficulties. With the growing demand for real-time train information and forecast, there is a pressing need to shift towards a data-driven, dynamic ETA prediction system that continuously adapts to actual train running conditions. Detailed\par\vspace{0.06cm}
\vspace{0.18cm}\noindent\textbf{\color{sihnavy}Description:}\par\vspace{0.06cm}
\noindent\small Indian Railways operates a vast network of passenger trains across diverse geographies, weather conditions, and traffic patterns. These coaching trains often face variability in journey times due to multiple real-world factors such as signal halts, congestion on busy routes, delays in preceding trains, temporary speed restrictions, unscheduled maintenance blocks, level crossing gates and operational bottlenecks.Despite this, ETA predictions at intermediate and destination stations are still often based on the train schedule, current delays and in-built recovery times, which lack accuracy and responsiveness.This limitation affects not just passengers but also impacts station planning, crew scheduling, platform allocation, cleaning operations, and feeder transport services. For long-distance trains with multi-day journeys, even a small deviation can cascade and lead to significant uncertainty. In an era where passengers expect real-time updates through mobile apps and station displays, inaccurate or outdated ETA predictions undermine service quality and trust. The challenge is to create a system that can dynamically forecast the ETA of trains at various points in their journey using real-time data feeds. These may include GPS-based location data, signal aspects, average sectional running times, weather conditions, historical delay patterns, and congestion levels on downstream tracks. The system must also be scalable to cover thousands of trains simultaneously and adaptable to the Indian Railways diverse operational zones. It should account for temporal and spatial variability in train performance and continuously refine its predictions using machine learning or statistical models. Such a system can serve as a foundation for better passenger communication, resource planning, and delay management.\par\vspace{0.06cm}
\vspace{0.18cm}\noindent\textbf{\color{sihnavy}Expected Solution:}\par\vspace{0.06cm}
\noindent\small The expected solution is a real-time ETA prediction system for coaching trains using data-driven models.It should integrate live train location data, operational parameters, historical delay trends, and network conditions to forecast arrival times at upcoming stations.The system must dynamically update ETAs in response to real-time events and delays. Machine learning or statistical forecasting techniques should be employed to improve accuracy over time. The solution should feature APIs for integration with mobile apps, station displays, and control room dashboards, ensuring that passengers and staff receive reliable, up-to-date information to support decision-making and planning.\par\vspace{0.06cm}

\end{tcolorbox}
\vspace{0.4cm}
\needspace{6\baselineskip}
\phantomsection
\addcontentsline{toc}{subsection}{[7.2/10] [SIH26073] AI/ML-Based Intelligent Anomaly Detection for Automatic Weather Stations (A...}

\begin{tcolorbox}[pscard, title={"\textbf{SIH26073} \quad $\vert$ \quad AI/ML-Based Intelligent Anomaly Detection for Automatic Weather Stations (A...}]
\textbf{\large AI/ML-Based Intelligent Anomaly Detection for Automatic Weather Stations (AWS)}\\[0.25cm]

\begin{tabular}{@{}lp{12.5cm}@{}}
\textbf{Difficulty Rating:} & \textcolor{diffhard}{\textbf{\Large 7.2 / 10}} \quad (\textbf{Hard (Intermediate-High)}) \\
\textbf{Problem ID:} & 26073 \quad $\vert$ \quad \textbf{Category:} \textbf{Software} \quad $\vert$ \quad \textbf{Theme:} \textbf{Disaster Management} \\
\textbf{Organization:} & Ministry of Earth Sciences (MoES) \\
\textbf{Department:} & India Meteorological Department \\
\textbf{Recommended Stack:} & \texttt{PyTorch, Scikit-Learn, XGBoost, Pandas, FastAPI, Celery / Redis} \\
\textbf{Key Challenges:} & \textit{\small Predictive ML \& anomaly detection models; Multi-source geospatial GIS \& real-time sensor fusion; Offline-first sync \& low-bandwidth field mobile apps; Disaster response planning \& resource management} \\
\end{tabular}

\vspace{0.2cm}
\hrule
\vspace{0.2cm}

\textbf{\color{sihnavy}Problem Statement Scope \& Requirements:}\\[0.15cm]
\noindent\hspace*{1.2em}\hangindent=2.4em\hangafter=1\textcolor{sihblue}{$\bullet$}~\small Title SkyGuard AI: Intelligent Real-Time Anomaly Detection System for Temperature, Pressure, and Humidity Sensors in Automatic Weather Stations\par\vspace{0.05cm}
\noindent\hspace*{1.2em}\hangindent=2.4em\hangafter=1\textcolor{sihblue}{$\bullet$}~\small Background Automatic Weather Stations (AWS) are critical components of modern meteorological observation networks. These stations continuously monitor atmospheric parameters and provide real-time data for weather forecasting, climate monitoring, disaster management, aviation, agriculture, and scientific research.However, AWS observations often contain anomalies caused by sensor malfunction,communication failures, calibration drift, power fluctuations, harsh environmental conditions, and data corruption.Erroneous observations can significantly impact weather forecasting accuracy and decision-making systems. Traditional threshold-based quality control methods are often insufficient for identifying complex or hidden anomalies in meteorological data streams.\par\vspace{0.05cm}
\noindent\hspace*{1.2em}\hangindent=2.4em\hangafter=1\textcolor{sihblue}{$\bullet$}~\small Problem Statement Develop an AI/ML-based intelligent anomaly detection system capable of automatically identifying abnormal, inconsistent, or faulty observations from Automatic Weather Stations in real time using only the following parameters:\par\vspace{0.05cm}
\noindent\hspace*{1.2em}\hangindent=2.4em\hangafter=1\textcolor{sihblue}{$\bullet$}~\small Temperature (Â$^\circ$C)\par\vspace{0.05cm}
\noindent\hspace*{1.2em}\hangindent=2.4em\hangafter=1\textcolor{sihblue}{$\bullet$}~\small Atmospheric Pressure (hPa)\par\vspace{0.05cm}
\noindent\hspace*{1.2em}\hangindent=2.4em\hangafter=1\textcolor{sihblue}{$\bullet$}~\small Relative Humidity (\%) The system should distinguish between genuine meteorological events and sensor/data anomalies while minimizing false alarms and enabling scalable deployment across large weather observation networks.\par\vspace{0.05cm}
\noindent\hspace*{1.2em}\hangindent=2.4em\hangafter=1\textcolor{sihblue}{$\bullet$}~\small Objectives\par\vspace{0.05cm}
\noindent\hspace*{1.2em}\hangindent=2.4em\hangafter=1\textcolor{sihblue}{$\bullet$}~\small Detect anomalies in real-time AWS data streams.\par\vspace{0.05cm}
\noindent\hspace*{1.2em}\hangindent=2.4em\hangafter=1\textcolor{sihblue}{$\bullet$}~\small Identify sensor faults, spikes, frozen values, and communication errors.\par\vspace{0.05cm}
\noindent\hspace*{1.2em}\hangindent=2.4em\hangafter=1\textcolor{sihblue}{$\bullet$}~\small Learn normal temporal and seasonal patterns of temperature, pressure, and humidity.\par\vspace{0.05cm}
\noindent\hspace*{1.2em}\hangindent=2.4em\hangafter=1\textcolor{sihblue}{$\bullet$}~\small Perform multivariate consistency analysis among atmospheric parameters.\par\vspace{0.05cm}
\noindent\hspace*{1.2em}\hangindent=2.4em\hangafter=1\textcolor{sihblue}{$\bullet$}~\small Provide confidence scores and explainable AI-based reasoning for detected anomalies.\par\vspace{0.05cm}
\noindent\hspace*{1.2em}\hangindent=2.4em\hangafter=1\textcolor{sihblue}{$\bullet$}~\small Predict possible sensor degradation and maintenance requirements.\par\vspace{0.05cm}
\noindent\hspace*{1.2em}\hangindent=2.4em\hangafter=1\textcolor{sihblue}{$\bullet$}~\small Optionally suggest corrected/imputed values for anomalous observations.\par\vspace{0.05cm}
\noindent\hspace*{1.2em}\hangindent=2.4em\hangafter=1\textcolor{sihblue}{$\bullet$}~\small Expected Inputs Participants may use historical AWS datasets, simulated anomalies, or streaming sensor data containing the following meteorological parameters: Parameter- Unit Temperature - Â$^\circ$C Atmospheric Pressure - hPa Relative Humidity - \%\par\vspace{0.05cm}
\noindent\hspace*{1.2em}\hangindent=2.4em\hangafter=1\textcolor{sihblue}{$\bullet$}~\small Expected Outputs\par\vspace{0.05cm}
\noindent\hspace*{1.2em}\hangindent=2.4em\hangafter=1\textcolor{sihblue}{$\bullet$}~\small Real-time anomaly alerts\par\vspace{0.05cm}
\noindent\hspace*{1.2em}\hangindent=2.4em\hangafter=1\textcolor{sihblue}{$\bullet$}~\small Severity and confidence scores\par\vspace{0.05cm}
\noindent\hspace*{1.2em}\hangindent=2.4em\hangafter=1\textcolor{sihblue}{$\bullet$}~\small Root-cause classification\par\vspace{0.05cm}
\noindent\hspace*{1.2em}\hangindent=2.4em\hangafter=1\textcolor{sihblue}{$\bullet$}~\small Visualization dashboard\par\vspace{0.05cm}
\noindent\hspace*{1.2em}\hangindent=2.4em\hangafter=1\textcolor{sihblue}{$\bullet$}~\small Sensor health status\par\vspace{0.05cm}
\noindent\hspace*{1.2em}\hangindent=2.4em\hangafter=1\textcolor{sihblue}{$\bullet$}~\small Corrected data estimation (optional)\par\vspace{0.05cm}
\noindent\hspace*{1.2em}\hangindent=2.4em\hangafter=1\textcolor{sihblue}{$\bullet$}~\small Suggested Technologies\par\vspace{0.05cm}
\noindent\hspace*{1.2em}\hangindent=2.4em\hangafter=1\textcolor{sihblue}{$\bullet$}~\small Explainable AI (SHAP/LIME) (Preferable)\par\vspace{0.05cm}
\noindent\hspace*{1.2em}\hangindent=2.4em\hangafter=1\textcolor{sihblue}{$\bullet$}~\small Edge AI for low-power deployment on ESP32\par\vspace{0.05cm}
\noindent\hspace*{1.2em}\hangindent=2.4em\hangafter=1\textcolor{sihblue}{$\bullet$}~\small Evaluation Criteria (To be evaluated in anomaly injected data) Criteria - Weightage Innovation \& Novelty - 25\% Detection Accuracy - 20\% Real-Time Capability - 15\% Explainability - 10\% Scalability - 10\% Practical Deployability - 10\% Visualization/UI - 5\% Energy Efficiency - 5\%\par\vspace{0.05cm}
\noindent\hspace*{1.2em}\hangindent=2.4em\hangafter=1\textcolor{sihblue}{$\bullet$}~\small Example Use Case An AWS suddenly reports a temperature of 55Â$^\circ$C with extremely high humidity and abnormal pressure variation while neighboring stations show normal conditions. The AI system should analyze temporal and spatial consistency, identify the reading as a probable sensor anomaly, generate an alert, and suggest corrective action.\par\vspace{0.05cm}
\noindent\hspace*{1.2em}\hangindent=2.4em\hangafter=1\textcolor{sihblue}{$\bullet$}~\small Grand Challenge Can AI build a self-aware and self-healing weather observation network capable of delivering trustworthy atmospheric data under all environmental conditions?\par\vspace{0.05cm}
\noindent\hspace*{1.2em}\hangindent=2.4em\hangafter=1\textcolor{sihblue}{$\bullet$}~\small Output: Fully executable code with example usage and a document explaining various use cases\par\vspace{0.05cm}

\end{tcolorbox}
\vspace{0.4cm}
\needspace{6\baselineskip}
\phantomsection
\addcontentsline{toc}{subsection}{[7.2/10] [SIH26192] Flash Flood Prediction System for Hilly Regions using Multi-Source Data The...}

\begin{tcolorbox}[pscard, title={"\textbf{SIH26192} \quad $\vert$ \quad Flash Flood Prediction System for Hilly Regions using Multi-Source Data The...}]
\textbf{\large Flash Flood Prediction System for Hilly Regions using Multi-Source Data Theme}\\[0.25cm]

\begin{tabular}{@{}lp{12.5cm}@{}}
\textbf{Difficulty Rating:} & \textcolor{diffhard}{\textbf{\Large 7.2 / 10}} \quad (\textbf{Hard (Intermediate-High)}) \\
\textbf{Problem ID:} & 26192 \quad $\vert$ \quad \textbf{Category:} \textbf{Software} \quad $\vert$ \quad \textbf{Theme:} \textbf{Disaster Management} \\
\textbf{Organization:} & Ministry of Home Affairs \\
\textbf{Department:} & National Disaster Response Force (NDRF), DM Division \\
\textbf{Recommended Stack:} & \texttt{PyTorch, Scikit-Learn, XGBoost, Pandas, GeoPandas, PostGIS} \\
\textbf{Key Challenges:} & \textit{\small Predictive ML \& anomaly detection models; Multi-source geospatial GIS \& real-time sensor fusion; Cloud web dashboard \& RESTful API microservices; Geotechnical vulnerability \& multi-hazard spatial zoning} \\
\end{tabular}

\vspace{0.2cm}
\hrule
\vspace{0.2cm}

\textbf{\color{sihnavy}Problem Statement Scope \& Requirements:}\\[0.15cm]
\noindent\hspace*{1.2em}\hangindent=2.4em\hangafter=1\textcolor{sihblue}{$\bullet$}~\small Background Hilly states in India are highly vulnerable to landslides and flash floods,which often occur with very short warning times. These sudden events result in significant loss of lives and property, and current early warning mechanisms are inadequate for hyper-local prediction and timely evacuation.\par\vspace{0.05cm}
\noindent\hspace*{1.2em}\hangindent=2.4em\hangafter=1\textcolor{sihblue}{$\bullet$}~\small Description The proposed initiative aims to develop a predictive system that integrates multiple data sources - rainfall data, soil moisture sensors,slope stability models, historical landslide inventories, and real-time IoT inputs. By combining these datasets, the system will generate hyper-local forecasts at the village or ward level, providing sufficient lead time for evacuation and risk mitigation.\par\vspace{0.05cm}
\noindent\hspace*{1.2em}\hangindent=2.4em\hangafter=1\textcolor{sihblue}{$\bullet$}~\small Expected Solution A comprehensive flash flood prediction system that Integrates rainfall,soil moisture, slope stability, and historical disaster data, utilizes IoT sensors for real-time monitoring, issues hyper-local early warnings at village/ward level, and provides actionable lead time for evacuation and disaster preparedness.\par\vspace{0.05cm}

\end{tcolorbox}
\vspace{0.4cm}
\needspace{6\baselineskip}
\phantomsection
\addcontentsline{toc}{subsection}{[7.1/10] [SIH26068] WeatherGPT: Conversational AI for Weather Forecasting, Alerts, and Climate ...}

\begin{tcolorbox}[pscard, title={"\textbf{SIH26068} \quad $\vert$ \quad WeatherGPT: Conversational AI for Weather Forecasting, Alerts, and Climate ...}]
\textbf{\large WeatherGPT: Conversational AI for Weather Forecasting, Alerts, and Climate Information}\\[0.25cm]

\begin{tabular}{@{}lp{12.5cm}@{}}
\textbf{Difficulty Rating:} & \textcolor{diffhard}{\textbf{\Large 7.1 / 10}} \quad (\textbf{Hard (Intermediate-High)}) \\
\textbf{Problem ID:} & 26068 \quad $\vert$ \quad \textbf{Category:} \textbf{Software} \quad $\vert$ \quad \textbf{Theme:} \textbf{Disaster Management} \\
\textbf{Organization:} & Ministry of Earth Sciences (MoES) \\
\textbf{Department:} & India Meteorological Department \\
\textbf{Recommended Stack:} & \texttt{PyTorch / TensorFlow, GDAL / Rasterio, Xarray, NetCDF4 / cfgrib, GeoPandas, PostGIS} \\
\textbf{Key Challenges:} & \textit{\small Physics-informed / Numerical simulation \& fluid dynamics coupling; Standard API feeds \& tabular time-series data; Cloud web dashboard \& RESTful API microservices; Disaster response planning \& resource management} \\
\end{tabular}

\vspace{0.2cm}
\hrule
\vspace{0.2cm}

\textbf{\color{sihnavy}Problem Statement Scope \& Requirements:}\\[0.15cm]
\noindent\hspace*{1.2em}\hangindent=2.4em\hangafter=1\textcolor{sihblue}{$\bullet$}~\small Background Weather information is often distributed through multiple portals, bulletins, satellite products, and forecast systems, making it difficult for common users, researchers, disaster managers, and government agencies to quickly obtain actionable insights. There is a need for an intelligent conversational platform that can provide real-time weather information, forecasts, warnings, climate analysis, and decision support in natural language.\par\vspace{0.05cm}
\noindent\hspace*{1.2em}\hangindent=2.4em\hangafter=1\textcolor{sihblue}{$\bullet$}~\small Objective Develop an AI-powered chatbot platform named WeatherGPT that integrates meteorological datasets, forecasting models, and disaster warning systems to provide accurate, contextual, and multilingual weather intelligence through conversational interfaces.\par\vspace{0.05cm}
\noindent\hspace*{1.2em}\hangindent=2.4em\hangafter=1\textcolor{sihblue}{$\bullet$}~\small Key Features\par\vspace{0.05cm}
\noindent\hangindent=1.8em\hangafter=1\textbf{1.}~\small Real-time weather information retrieval.\par\vspace{0.06cm}
\noindent\hangindent=1.8em\hangafter=1\textbf{2.}~\small Natural language querying for weather forecasts.\par\vspace{0.06cm}
\noindent\hangindent=1.8em\hangafter=1\textbf{3.}~\small Integration with numerical weather prediction (NWP) models such as GFS/WRF.\par\vspace{0.06cm}
\noindent\hangindent=1.8em\hangafter=1\textbf{4.}~\small Extreme weather alerts and early warning dissemination.\par\vspace{0.06cm}
\noindent\hangindent=1.8em\hangafter=1\textbf{5.}~\small Location-based forecasting and advisory generation.\par\vspace{0.06cm}
\noindent\hangindent=1.8em\hangafter=1\textbf{6.}~\small Multilingual support for Indian languages.\par\vspace{0.06cm}
\noindent\hangindent=1.8em\hangafter=1\textbf{7.}~\small Climate trend and historical weather analysis.\par\vspace{0.06cm}
\noindent\hangindent=1.8em\hangafter=1\textbf{8.}~\small Voice-enabled interaction for rural accessibility.\par\vspace{0.06cm}
\noindent\hspace*{1.2em}\hangindent=2.4em\hangafter=1\textcolor{sihblue}{$\bullet$}~\small Expected Solution Participants should develop:\par\vspace{0.05cm}
\noindent\hspace*{1.2em}\hangindent=2.4em\hangafter=1\textcolor{sihblue}{$\bullet$}~\small A mobile-based conversational AI platform.\par\vspace{0.05cm}
\noindent\hspace*{1.2em}\hangindent=2.4em\hangafter=1\textcolor{sihblue}{$\bullet$}~\small Backend integration with meteorological databases, website and APIs.\par\vspace{0.05cm}
\noindent\hspace*{1.2em}\hangindent=2.4em\hangafter=1\textcolor{sihblue}{$\bullet$}~\small AI/LLM-based query understanding engine.\par\vspace{0.05cm}
\noindent\hspace*{1.2em}\hangindent=2.4em\hangafter=1\textcolor{sihblue}{$\bullet$}~\small Scalable architecture supporting real-time data ingestion.\par\vspace{0.05cm}
\noindent\hspace*{1.2em}\hangindent=2.4em\hangafter=1\textcolor{sihblue}{$\bullet$}~\small Suggested Technology Stack\par\vspace{0.05cm}
\noindent\hspace*{1.2em}\hangindent=2.4em\hangafter=1\textcolor{sihblue}{$\bullet$}~\small Python / FastAPI / Node.js\par\vspace{0.05cm}
\noindent\hspace*{1.2em}\hangindent=2.4em\hangafter=1\textcolor{sihblue}{$\bullet$}~\small MQTT / WIS2.0 / WebSocket\par\vspace{0.05cm}
\noindent\hspace*{1.2em}\hangindent=2.4em\hangafter=1\textcolor{sihblue}{$\bullet$}~\small LLMs (OpenAI, Llama, Gemini, etc.)\par\vspace{0.05cm}
\noindent\hspace*{1.2em}\hangindent=2.4em\hangafter=1\textcolor{sihblue}{$\bullet$}~\small GIS tools and weather APIs\par\vspace{0.05cm}
\noindent\hspace*{1.2em}\hangindent=2.4em\hangafter=1\textcolor{sihblue}{$\bullet$}~\small PostgreSQL / MongoDB\par\vspace{0.05cm}
\noindent\hspace*{1.2em}\hangindent=2.4em\hangafter=1\textcolor{sihblue}{$\bullet$}~\small Docker / Kubernetes\par\vspace{0.05cm}
\noindent\hspace*{1.2em}\hangindent=2.4em\hangafter=1\textcolor{sihblue}{$\bullet$}~\small Expected Outcomes\par\vspace{0.05cm}
\noindent\hspace*{1.2em}\hangindent=2.4em\hangafter=1\textcolor{sihblue}{$\bullet$}~\small Faster dissemination of weather information.\par\vspace{0.05cm}
\noindent\hspace*{1.2em}\hangindent=2.4em\hangafter=1\textcolor{sihblue}{$\bullet$}~\small Improved public accessibility to forecasts.\par\vspace{0.05cm}
\noindent\hspace*{1.2em}\hangindent=2.4em\hangafter=1\textcolor{sihblue}{$\bullet$}~\small Better disaster preparedness and response.\par\vspace{0.05cm}
\noindent\hspace*{1.2em}\hangindent=2.4em\hangafter=1\textcolor{sihblue}{$\bullet$}~\small Intelligent weather decision-support system for agriculture, aviation, marine, and urban planning.\par\vspace{0.05cm}
\noindent\hspace*{1.2em}\hangindent=2.4em\hangafter=1\textcolor{sihblue}{$\bullet$}~\small Possible Use Cases\par\vspace{0.05cm}
\noindent\hspace*{1.2em}\hangindent=2.4em\hangafter=1\textcolor{sihblue}{$\bullet$}~\small Farmers seeking crop-weather advisories.\par\vspace{0.05cm}
\noindent\hspace*{1.2em}\hangindent=2.4em\hangafter=1\textcolor{sihblue}{$\bullet$}~\small Aviation weather briefing.\par\vspace{0.05cm}
\noindent\hspace*{1.2em}\hangindent=2.4em\hangafter=1\textcolor{sihblue}{$\bullet$}~\small Flood/cyclone warning dissemination.\par\vspace{0.05cm}
\noindent\hspace*{1.2em}\hangindent=2.4em\hangafter=1\textcolor{sihblue}{$\bullet$}~\small Smart city weather monitoring.\par\vspace{0.05cm}
\noindent\hspace*{1.2em}\hangindent=2.4em\hangafter=1\textcolor{sihblue}{$\bullet$}~\small Climate analytics for researchers.\par\vspace{0.05cm}
\noindent\hspace*{1.2em}\hangindent=2.4em\hangafter=1\textcolor{sihblue}{$\bullet$}~\small Evaluation Parameters\par\vspace{0.05cm}
\noindent\hspace*{1.2em}\hangindent=2.4em\hangafter=1\textcolor{sihblue}{$\bullet$}~\small Accuracy and relevance.\par\vspace{0.05cm}
\noindent\hspace*{1.2em}\hangindent=2.4em\hangafter=1\textcolor{sihblue}{$\bullet$}~\small Response latency.\par\vspace{0.05cm}
\noindent\hspace*{1.2em}\hangindent=2.4em\hangafter=1\textcolor{sihblue}{$\bullet$}~\small Multilingual capability.\par\vspace{0.05cm}
\noindent\hspace*{1.2em}\hangindent=2.4em\hangafter=1\textcolor{sihblue}{$\bullet$}~\small User interface and accessibility.\par\vspace{0.05cm}
\noindent\hspace*{1.2em}\hangindent=2.4em\hangafter=1\textcolor{sihblue}{$\bullet$}~\small Scalability and innovation.\par\vspace{0.05cm}
\noindent\hspace*{1.2em}\hangindent=2.4em\hangafter=1\textcolor{sihblue}{$\bullet$}~\small Integration with real-time meteorological systems.\par\vspace{0.05cm}
\noindent\hspace*{1.2em}\hangindent=2.4em\hangafter=1\textcolor{sihblue}{$\bullet$}~\small Voice-enabled interaction for rural accessibility\par\vspace{0.05cm}

\end{tcolorbox}
\vspace{0.4cm}
\needspace{6\baselineskip}
\phantomsection
\addcontentsline{toc}{subsection}{[7.1/10] [SIH26079] AI-Based Forecast Bust Detection for Medium-Range Weather Forecasts}

\begin{tcolorbox}[pscard, title={"\textbf{SIH26079} \quad $\vert$ \quad AI-Based Forecast Bust Detection for Medium-Range Weather Forecasts}]
\textbf{\large AI-Based Forecast Bust Detection for Medium-Range Weather Forecasts}\\[0.25cm]

\begin{tabular}{@{}lp{12.5cm}@{}}
\textbf{Difficulty Rating:} & \textcolor{diffhard}{\textbf{\Large 7.1 / 10}} \quad (\textbf{Hard (Intermediate-High)}) \\
\textbf{Problem ID:} & 26079 \quad $\vert$ \quad \textbf{Category:} \textbf{Software} \quad $\vert$ \quad \textbf{Theme:} \textbf{Disaster Management} \\
\textbf{Organization:} & Ministry of Earth Sciences (MoES) \\
\textbf{Department:} & National Centre for Medium Range Weather Forecasting (NCMRWF) \\
\textbf{Recommended Stack:} & \texttt{PyTorch, Scikit-Learn, XGBoost, Pandas, GeoPandas, PostGIS} \\
\textbf{Key Challenges:} & \textit{\small Physics-informed / Numerical simulation \& fluid dynamics coupling; Standard API feeds \& tabular time-series data; Cloud web dashboard \& RESTful API microservices; Disaster response planning \& resource management} \\
\end{tabular}

\vspace{0.2cm}
\hrule
\vspace{0.2cm}

\textbf{\color{sihnavy}Problem Statement Scope \& Requirements:}\\[0.15cm]
\noindent\hspace*{1.2em}\hangindent=2.4em\hangafter=1\textcolor{sihblue}{$\bullet$}~\small Problem Statement Medium-range weather forecasts sometimes show large errors during rapidly evolving systems such as monsoon depressions, heavy rainfall events, western disturbances, cyclones, heat waves and break/active monsoon phases. Such forecast failures, or 'forecast busts', can affect operational decision-making.\par\vspace{0.05cm}
\noindent\hspace*{1.2em}\hangindent=2.4em\hangafter=1\textcolor{sihblue}{$\bullet$}~\small Challenge The challenge is to develop an AI/ML-based system that can identify regions and lead times where the forecast is likely to have high uncertainty or large error. The system should compare current NWP forecast patterns with historical forecast error behaviour and provide a forecast confidence indicator. Expected Outcome - Description Forecast confidence map - Region-wise confidence for Day 1 to Day 10 forecasts Forecast bust probability - Probability of large forecast error over different regions Error-prone area detection - Identification of areas where model forecast may be unreliable Explainable output - Key meteorological reasons for low confidence Prototype dashboard/API - Simple interface for operational use\par\vspace{0.05cm}

\end{tcolorbox}
\vspace{0.4cm}
\needspace{6\baselineskip}
\phantomsection
\addcontentsline{toc}{subsection}{[7.1/10] [SIH26080] Regime-Aware AI Post-Processing of Monsoon Rainfall Forecasts}

\begin{tcolorbox}[pscard, title={"\textbf{SIH26080} \quad $\vert$ \quad Regime-Aware AI Post-Processing of Monsoon Rainfall Forecasts}]
\textbf{\large Regime-Aware AI Post-Processing of Monsoon Rainfall Forecasts}\\[0.25cm]

\begin{tabular}{@{}lp{12.5cm}@{}}
\textbf{Difficulty Rating:} & \textcolor{diffhard}{\textbf{\Large 7.1 / 10}} \quad (\textbf{Hard (Intermediate-High)}) \\
\textbf{Problem ID:} & 26080 \quad $\vert$ \quad \textbf{Category:} \textbf{Software} \quad $\vert$ \quad \textbf{Theme:} \textbf{Disaster Management} \\
\textbf{Organization:} & Ministry of Earth Sciences (MoES) \\
\textbf{Department:} & National Centre for Medium Range Weather Forecasting (NCMRWF) \\
\textbf{Recommended Stack:} & \texttt{PyTorch, Scikit-Learn, XGBoost, Pandas, GeoPandas, PostGIS} \\
\textbf{Key Challenges:} & \textit{\small Physics-informed / Numerical simulation \& fluid dynamics coupling; Standard API feeds \& tabular time-series data; Cloud web dashboard \& RESTful API microservices; Disaster response planning \& resource management} \\
\end{tabular}

\vspace{0.2cm}
\hrule
\vspace{0.2cm}

\textbf{\color{sihnavy}Problem Statement Scope \& Requirements:}\\[0.15cm]
\noindent\hspace*{1.2em}\hangindent=2.4em\hangafter=1\textcolor{sihblue}{$\bullet$}~\small Problem Statement Rainfall forecast errors over India vary with weather regimes such as active monsoon, break monsoon, monsoon lows/depressions, orographic rainfall, coastal rainfall and western disturbances. A single bias-correction method may not work equally well in all situations.The challenge is to build an AI/ML-based rainfall post-processing system that first identifies the prevailing weather regime and then applies suitable correction to the raw NWP rainfall forecast.The aim is to improve district/grid-level rainfall forecasts, especially for heavy and very heavy rainfall events.\par\vspace{0.05cm}
\noindent\hspace*{1.2em}\hangindent=2.4em\hangafter=1\textcolor{sihblue}{$\bullet$}~\small Expected Outcome Expected Outcome -\par\vspace{0.05cm}
\vspace{0.18cm}\noindent\textbf{\color{sihnavy}Description:}\par\vspace{0.06cm}
\noindent\small Weather regime classifier - Classification of active, break, depression,coastal/orographic rainfall regimes Bias-corrected rainfall forecast - Improved rainfall forecast compared to raw NWP output Heavy rainfall probability - Probability of rainfall exceeding operational thresholds District-level rainfall product - User-friendly rainfall forecast table/map Verification report - Skill comparison using RMSE, ETS, CSI, POD, FAR and FSS\par\vspace{0.06cm}

\end{tcolorbox}
\vspace{0.4cm}
\needspace{6\baselineskip}
\phantomsection
\addcontentsline{toc}{subsection}{[7.1/10] [SIH26082] Air PollutionWeather Coupled Forecasting System (Delhi NCR Focus)}

\begin{tcolorbox}[pscard, title={"\textbf{SIH26082} \quad $\vert$ \quad Air PollutionWeather Coupled Forecasting System (Delhi NCR Focus)}]
\textbf{\large Air PollutionWeather Coupled Forecasting System (Delhi NCR Focus)}\\[0.25cm]

\begin{tabular}{@{}lp{12.5cm}@{}}
\textbf{Difficulty Rating:} & \textcolor{diffhard}{\textbf{\Large 7.1 / 10}} \quad (\textbf{Hard (Intermediate-High)}) \\
\textbf{Problem ID:} & 26082 \quad $\vert$ \quad \textbf{Category:} \textbf{Software} \quad $\vert$ \quad \textbf{Theme:} \textbf{Disaster Management} \\
\textbf{Organization:} & Ministry of Earth Sciences (MoES) \\
\textbf{Department:} & National Centre for Medium Range Weather Forecasting (NCMRWF) \\
\textbf{Recommended Stack:} & \texttt{PyTorch, Scikit-Learn, XGBoost, Pandas, FastAPI, Celery / Redis} \\
\textbf{Key Challenges:} & \textit{\small Physics-informed / Numerical simulation \& fluid dynamics coupling; Standard API feeds \& tabular time-series data; Cloud web dashboard \& RESTful API microservices; Disaster response planning \& resource management} \\
\end{tabular}

\vspace{0.2cm}
\hrule
\vspace{0.2cm}

\textbf{\color{sihnavy}Problem Statement Scope \& Requirements:}\\[0.15cm]
\noindent\small Traditional Air Quality Index (AQI) forecasting models typically treat meteorology and pollution dispersion as separate entities. However, in highly polluted urban landscapes like Delhi NCR, there is a critical, dynamic feedback loop between the weather and pollutants. During peak pollution seasons (such as the winter stubble-burning period), atmospheric inversion layers trap particulate matter close to the ground. Conversely, dense concentrations of aerosols (PM2.5) block sunlight,altering local temperatures, wind patterns, and planetary boundary layer (PBL) heights. Ignoring these coupled meteorological-chemical feedback loops leads to significant inaccuracies in standard AQI predictions. To achieve high-accuracy, actionable insights, there is an urgent need for an integrated system that simulates real-time interactions between atmospheric physics and chemical transport. The challenge is to build a high-resolution, coupled forecasting system specifically tailored for Delhi NCR that predicts AQI for the next 72 hours. The solution must leverage advanced weather-chemistry models (such as WRF-Chem or similar open-source coupled frameworks) to dynamically interlink meteorology with pollution dispersion(specifically PM2.5 and Ground-level Ozone). A core focus should be accurately modeling the impact of atmospheric inversion on external pollution spikes, such as regional stubble burning,and how those trapped pollutants subsequently alter local weather conditions.Implement a workflow that handles two-way feedback between meteorology (temperature, wind,PBL height) and chemistry (PM2.5, PM10, O3,NOx). A user-friendly, real-time dashboard displaying high-resolution AQI forecasts for Delhi NCR with a 72-hour outlook. Features that explicitly track atmospheric inversion strength and predict how stubble-burning plumes will disperse under prevailing weather conditions.\par\vspace{0.06cm}

\end{tcolorbox}
\vspace{0.4cm}
\needspace{6\baselineskip}
\phantomsection
\addcontentsline{toc}{subsection}{[6.5/10] [SIH26015] Application of Geospatial Techniques for visualization and analysis to inte...}

\begin{tcolorbox}[pscard, title={"\textbf{SIH26015} \quad $\vert$ \quad Application of Geospatial Techniques for visualization and analysis to inte...}]
\textbf{\large Application of Geospatial Techniques for visualization and analysis to interpret Geo-Coded lmages to enhance watershed Development Outcomes.}\\[0.25cm]

\begin{tabular}{@{}lp{12.5cm}@{}}
\textbf{Difficulty Rating:} & \textcolor{diffhard}{\textbf{\Large 6.5 / 10}} \quad (\textbf{Hard (Intermediate-High)}) \\
\textbf{Problem ID:} & 26015 \quad $\vert$ \quad \textbf{Category:} \textbf{Software} \quad $\vert$ \quad \textbf{Theme:} \textbf{Disaster Management} \\
\textbf{Organization:} & Ministry of Rural Development \\
\textbf{Department:} & Dept of land resources (DoLR) \\
\textbf{Recommended Stack:} & \texttt{PyTorch / TensorFlow, GDAL / Rasterio, Xarray, NetCDF4 / cfgrib, GeoPandas, PostGIS} \\
\textbf{Key Challenges:} & \textit{\small Data processing \& rule-based decision engine; High-volume Doppler radar \& satellite raster/GRIB2 ingestion; Cloud web dashboard \& RESTful API microservices; Disaster response planning \& resource management} \\
\end{tabular}

\vspace{0.2cm}
\hrule
\vspace{0.2cm}

\textbf{\color{sihnavy}Problem Statement Scope \& Requirements:}\\[0.15cm]
\vspace{0.18cm}\noindent\textbf{\color{sihnavy}Background:}~\small Watershed development plays a vital role in sustainable management of land, water, and natural resources, particularly in rural and semi-arid regions of lndia. Effective watershed planning and monitoring require accurate spatial information on land use, drainage patterns, vegetation cover, soil moisture, water bodies, and changes occurring over time. Traditional monitoring approaches often rely on field surveys and manual reporting, which are time-consuming, resource-intensive, and limited in spatial coverage. ln recent years, advancements in Geographic lnformation systems (GlS), Remote Sensing (RS), and geospatial technologies have created new opportunities for scientific watershed assessment and evidence-based decision- making.Geo-coded images, integrated with satellite-based spatial datasets, provide location- specific visual information that can significantly improve watershed monitoring and interpretation. The availability of 30 m spatial resolution satellite data through the SRISHTI-DRISHTI platform offers a valuable opportunity to develop analytical frameworks for visualization, interpretation, and assessment of watershed characteristics. Through thematic mapping and spatial analysis, geo-coded images can support identification of land degradation, water conservation structures, vegetation changes, drainage conditions, and other watershed-related parameters. The proposed study focuses on the application of geospatial techniques for visualization and analysis to interpret geo-coded images for enhancing watershed development outcomes. The study aims to develop a systematic and scalable approach for image-based watershed analysis using GIS and remote sensing tools. By integrating geo-coded imagery with satellite datasets from the SRISHTI-DRISHTI platform, the research seeks to improve planning, monitoring, interpretation, and scientific assessment of watershed interventions in a focused and cost-effective manner. Description of the Study: Despite significant investments in watershed development programs, effective monitoring and interpretation of watershed activities remain major challenges. Existing assessment methods are often fragmented, dependent on manual observations, and lack spatial integration. Many watershed projects face difficulties in accurately visualizing field conditions, tracking spatial changes, identifying intervention impacts, and generating reliable evidence for decision-making. Although geo-coded images are increasingly being collected during watershed implementation and monitoring, their analytical utilization remains limited. ln many cases, geo-tagged photographs are used only for documentation purposes rather than for integrated spatial analysis and interpretation. There is insufficient use of advanced GIS and remote sensing techniques to systematically visualize, analyse, and interpret these geo-coded datasets in relation to watershed characteristics and satellite-derived information. Furthermore, the absence of standardized visualization frameworks restricts the ability of planners and administrators to derive actionable insights from geo-coded imagery. Challenges also exist in integrating field-level geo-coded images with satellite data, thematic layers, and watershed boundaries for meaningful analysis. Limited technical approaches for image interpretation reduce the effectiveness of watershed monitoring systems and hinder scientific evaluation of land and water resource interventions. The SRISHTI-DRISHTI platform provides an opportunity to address these challenges by offering consistent 30 m spatial resolution satellite data that can support integrated geospatial analysis. However, there is a need to develop specialized methodologies and visualization techniques that can effectively interpret geo-coded images and generate meaningful watershed insights. Therefore, the study aims to bridge this gap through a focused analytical framework combining GlS, remote sensing, thematic mapping, and geo-coded image interpretation for enhanced watershed development outcomes. Scope of the Study: Table to be Added here Expected Solutions: The proposed study is expected to provide a structured geospatial framework for visualization and interpretation of geo-coded images in watershed development programs. The key expected solutions include: a)Development of an lntegrated Geospatial Visualization Framework: Thes tudy will develop a GIS and remote sensing-based framework for integrating geo-coded images with satellite data from the SRISHTI-DRISHTI platform to support watershed analysis and monitoring. b) lmproved Geo-Coded lmage lnterpretation: Advanced spatial interpretation techniques will help convert geo-coded imagery into meaningful analytical information related to land use, Vegetation status, water resources, watershed interventions, and environmental changes. c) Generation of Thematic Maps and Visualization Products: The study will visualize outputs such as land use maps,drainage maps, vegetation maps, watershed intervention maps, and spatial change detection products for better interpretation and planning. d) Enhanced Watershed Monitoring and Assessment: lntegration of geo-coded images with satellite-based datasets will improve the accuracy and efficiency of monitoring watershed activities and assessing development outcomes. e) Scientific Support for Decision-Making: The proposed framework will support evidence-based planning and policy decisions by providing spatially validated and visually interpretable watershed information. f) Scalable and Cost-Effective Monitoring Approach: The methodology will offer a focused, scalable, and economically viable approach for watershed monitoring that can be replicated across different regions and watershed programs. g) Strengthening Use of the SRISHTI-DRISHTI Platform: The study will enhance the practical utilization of the SRISHTI-DRISHTI platform as a dedicated source for satellite-based geospatial analysis and watershed interpretation.\par\vspace{0.06cm}

\end{tcolorbox}
\vspace{0.4cm}
\needspace{6\baselineskip}
\phantomsection
\addcontentsline{toc}{subsection}{[6.5/10] [SIH26043] A digital platform to crowdsource societal challenges and facilitate collab...}

\begin{tcolorbox}[pscard, title={"\textbf{SIH26043} \quad $\vert$ \quad A digital platform to crowdsource societal challenges and facilitate collab...}]
\textbf{\large A digital platform to crowdsource societal challenges and facilitate collaborative problem solving through universities and industry partnerships}\\[0.25cm]

\begin{tabular}{@{}lp{12.5cm}@{}}
\textbf{Difficulty Rating:} & \textcolor{diffhard}{\textbf{\Large 6.5 / 10}} \quad (\textbf{Hard (Intermediate-High)}) \\
\textbf{Problem ID:} & 26043 \quad $\vert$ \quad \textbf{Category:} \textbf{Software} \quad $\vert$ \quad \textbf{Theme:} \textbf{Disaster Management} \\
\textbf{Organization:} & Governmcnt of Jharkhand \\
\textbf{Department:} & Department of Higher \& Technical Education \\
\textbf{Recommended Stack:} & \texttt{PyTorch, Scikit-Learn, XGBoost, Pandas, FastAPI, Celery / Redis} \\
\textbf{Key Challenges:} & \textit{\small Data processing \& rule-based decision engine; High-volume Doppler radar \& satellite raster/GRIB2 ingestion; Cloud web dashboard \& RESTful API microservices; Disaster response planning \& resource management} \\
\end{tabular}

\vspace{0.2cm}
\hrule
\vspace{0.2cm}

\textbf{\color{sihnavy}Problem Statement Scope \& Requirements:}\\[0.15cm]
\vspace{0.18cm}\noindent\textbf{\color{sihnavy}Background:}~\small Communities across Jharkhand encounter numerous local challenges related to education,healthcare, agriculture, water management, sanitation, environment, rural livelihoods,accessibility, urban infrastructure, and public service delivery. While citizens are often the first to identify these issues, there is currently no structured mechanism through which they can submit such problems for systematic evaluation and innovation-driven resolution.At the same time, Higher Education lnstitutions (HEIs) possess significant academic expertise,research capabilities, and a large pool of students capable of developing practical solutions.Industries and start-ups also have technical expertise, financial resources, and implementation capabilities that can complement academic research. However, collaboration among citizens,universities, and industry remains largely fragmented and project-specific.The National Education Policy (NEP) 2020 emphasizes experiential learning, multidisciplinary research, innovation, industry collaboration, and community engagement. Establishing a technology-enabled platform that connects societal challenges with academic institutions and industry partners can foster demand-driven innovation while enabling students and researchers to work on real-world problems that create measurable social impact.\par\vspace{0.06cm}
\vspace{0.18cm}\noindent\textbf{\color{sihnavy}Description:}\par\vspace{0.06cm}
\noindent\small Every year, citizens across Jharkhand identiff thousands of local issues that require innovative technological or process-based solutions. These challenges often remain unresolved due to the absence of a centralized platform that enables problem collection, categorization, expert evaluation, institutional assignment, and industry collaboration.There is a need to develop a digital platform capable of:\par\vspace{0.06cm}
\noindent\hspace*{1.2em}\hangindent=2.4em\hangafter=1\textcolor{sihblue}{$\bullet$}~\small Allowing citizens, community organizations, local bodies, and government agencies to submit societal challenges thiough an intuitive web and mobile interface, supported by photographs, videos, location details, and relevant documents.\par\vspace{0.05cm}
\noindent\hspace*{1.2em}\hangindent=2.4em\hangafter=1\textcolor{sihblue}{$\bullet$}~\small Automatically categorizing submitted problems based on thematic domains such as education, agriculture, healthcare, water resources, environment, energy, urban development, accessibility, public administration, and rural livelihoods using Al-enabled classification techniques.\par\vspace{0.05cm}
\noindent\hspace*{1.2em}\hangindent=2.4em\hangafter=1\textcolor{sihblue}{$\bullet$}~\small Routing validated problem statements to appropriate universities based on their academic disciplines, research expertise, innovation centres, incubation facilities, and faculty specialization.\par\vspace{0.05cm}
\noindent\hspace*{1.2em}\hangindent=2.4em\hangafter=1\textcolor{sihblue}{$\bullet$}~\small Enabling universities to evaluate submitted challenges, constitute multidisciplinary student and faculty teams, and prepare solution proposals or research projects.\par\vspace{0.05cm}
\noindent\hspace*{1.2em}\hangindent=2.4em\hangafter=1\textcolor{sihblue}{$\bullet$}~\small Facilitating collaboration between universities and industry partners, startups, MSMEs,CSR organizations, research laboratories, and innovation ecosystems for mentorship,funding, prototyping, testing and deployment of solutions.\par\vspace{0.05cm}
\noindent\hspace*{1.2em}\hangindent=2.4em\hangafter=1\textcolor{sihblue}{$\bullet$}~\small Providing workflow management for problem review, institutional allocation, project monitoring, stakeholder communication, milestone tracking, and solution validation.\par\vspace{0.05cm}
\noindent\hspace*{1.2em}\hangindent=2.4em\hangafter=1\textcolor{sihblue}{$\bullet$}~\small Generating dashboards and analytics for govemment departments to monitor the number of challenges received, domain-wise distribution, institutional participation, industry engagement, project progress, and measurable social outcomes. The platform should support a transparent and scalable innovation ecosystem that transforms community-driven challenges into research, innovation, entrepreneurship, and deployable solutions.\par\vspace{0.05cm}
\vspace{0.18cm}\noindent\textbf{\color{sihnavy}Expected Solution:}\par\vspace{0.06cm}
\noindent\small A comprehensive Societal Innovation Collaboration Portal comprising the following components:\par\vspace{0.06cm}
\noindent\hspace*{1.2em}\hangindent=2.4em\hangafter=1\textcolor{sihblue}{$\bullet$}~\small A citizen engagement module enabling individuals, community groups, Panchayati Raj Institutions, Urban Local Bodies, and government departments to submit societal challenges with multimedia evidence, geographical location, and supporting information.\par\vspace{0.05cm}
\noindent\hspace*{1.2em}\hangindent=2.4em\hangafter=1\textcolor{sihblue}{$\bullet$}~\small An Al-enabled problem management module capable of automatically categorizing, prioritizing, deduplication, and routing validated challenges to appropriate universities based on subject expertise and institutional capabilities.\par\vspace{0.05cm}
\noindent\hspace*{1.2em}\hangindent=2.4em\hangafter=1\textcolor{sihblue}{$\bullet$}~\small A university collaboration module allowing Higher Education Institutions to review assigned challenges, form multidisciplinary project teams, assign faculty mentors, manage project workflows, and submit solution proposals.\par\vspace{0.05cm}
\noindent\hspace*{1.2em}\hangindent=2.4em\hangafter=1\textcolor{sihblue}{$\bullet$}~\small An industry partnership module facilitating participation by industries, startups, MSMEs,CSR organizations, research institutions, and innovation hubs for mentoring, co-development, funding, prototyping, pilot implementation, and technology transfer.\par\vspace{0.05cm}
\noindent\hspace*{1.2em}\hangindent=2.4em\hangafter=1\textcolor{sihblue}{$\bullet$}~\small A project lifecycle management system for monitoring milestones, deliverables, approvals, documentation, testing outcomes, intellectual property generation, and implementation status.\par\vspace{0.05cm}
\noindent\hspace*{1.2em}\hangindent=2.4em\hangafter=1\textcolor{sihblue}{$\bullet$}~\small A visual analytics dashboard providing real-time insights on challenge submissions, university participation, industry collaborations, thematic trends, project completion rates,innovation outcomes, patents, startups created, and community impact across districts and sectors.\par\vspace{0.05cm}
\noindent\hspace*{1.2em}\hangindent=2.4em\hangafter=1\textcolor{sihblue}{$\bullet$}~\small A notification and communication system enabling seamless interaction among citizens,universities, industry partners, mentors, and government departments throughout the project lifecycle.\par\vspace{0.05cm}

\end{tcolorbox}
\vspace{0.4cm}
\needspace{6\baselineskip}
\phantomsection
\addcontentsline{toc}{subsection}{[6.4/10] [SIH26069] National Weather Big Data Analytics Platform}

\begin{tcolorbox}[pscard, title={"\textbf{SIH26069} \quad $\vert$ \quad National Weather Big Data Analytics Platform}]
\textbf{\large National Weather Big Data Analytics Platform}\\[0.25cm]

\begin{tabular}{@{}lp{12.5cm}@{}}
\textbf{Difficulty Rating:} & \textcolor{diffmod}{\textbf{\Large 6.4 / 10}} \quad (\textbf{Moderate (Intermediate)}) \\
\textbf{Problem ID:} & 26069 \quad $\vert$ \quad \textbf{Category:} \textbf{Software} \quad $\vert$ \quad \textbf{Theme:} \textbf{Disaster Management} \\
\textbf{Organization:} & Ministry of Earth Sciences (MoES) \\
\textbf{Department:} & India Meteorological Department \\
\textbf{Recommended Stack:} & \texttt{PyTorch, Scikit-Learn, XGBoost, Pandas, GeoPandas, PostGIS} \\
\textbf{Key Challenges:} & \textit{\small Predictive ML \& anomaly detection models; Standard API feeds \& tabular time-series data; Cloud web dashboard \& RESTful API microservices; Disaster response planning \& resource management} \\
\end{tabular}

\vspace{0.2cm}
\hrule
\vspace{0.2cm}

\textbf{\color{sihnavy}Problem Statement Scope \& Requirements:}\\[0.15cm]
\noindent\small Design and develop a scalable National Weather Big Data Analytics Platform capable of collecting and processing real-time weather-related information for India from multiple internet-based sources including social media platforms, public datasets, websites, APIs, and citizen reports. The platform should automatically collect weather related posts and information tagged with \#IMD and other relevant weather hashtags, along with metadata such as date \& time, city, state, GPS location, photos, videos, and event category, and store the information in a centralized database. The system should leverage big data technologies and open-source tools to support large-scale real-time data ingestion, processing, storage, and visualization. Participants are encouraged to use machine learning and AI-based techniques to identify fake or misleading reports, verify untrusted sources, remove duplicate entries, and automatically categorize weather events such as rainfall, thunderstorms, flooding, heatwaves, fog, dust storms, and strong winds. Develop a web-based dashboard and Admin Panel for monitoring and analysing collected data with features including:\par\vspace{0.06cm}
\noindent\hspace*{1.2em}\hangindent=2.4em\hangafter=1\textcolor{sihblue}{$\bullet$}~\small Date-wise filtering\par\vspace{0.05cm}
\noindent\hspace*{1.2em}\hangindent=2.4em\hangafter=1\textcolor{sihblue}{$\bullet$}~\small Event-wise filtering\par\vspace{0.05cm}
\noindent\hspace*{1.2em}\hangindent=2.4em\hangafter=1\textcolor{sihblue}{$\bullet$}~\small Location-wise filtering\par\vspace{0.05cm}
\noindent\hspace*{1.2em}\hangindent=2.4em\hangafter=1\textcolor{sihblue}{$\bullet$}~\small Verification status tracking\par\vspace{0.05cm}
\noindent\hspace*{1.2em}\hangindent=2.4em\hangafter=1\textcolor{sihblue}{$\bullet$}~\small Real-time visualization and analytics\par\vspace{0.05cm}

\end{tcolorbox}
\vspace{0.4cm}
\needspace{6\baselineskip}
\phantomsection
\addcontentsline{toc}{subsection}{[5.9/10] [SIH26191] Intelligent Identification of Hazard-Based Red Zones, Carrying Capacity Ass...}

\begin{tcolorbox}[pscard, title={"\textbf{SIH26191} \quad $\vert$ \quad Intelligent Identification of Hazard-Based Red Zones, Carrying Capacity Ass...}]
\textbf{\large Intelligent Identification of Hazard-Based Red Zones, Carrying Capacity Assessment, and Immediate Relocation Needs for Vulnerable Habitations}\\[0.25cm]

\begin{tabular}{@{}lp{12.5cm}@{}}
\textbf{Difficulty Rating:} & \textcolor{diffmod}{\textbf{\Large 5.9 / 10}} \quad (\textbf{Moderate (Intermediate)}) \\
\textbf{Problem ID:} & 26191 \quad $\vert$ \quad \textbf{Category:} \textbf{Software} \quad $\vert$ \quad \textbf{Theme:} \textbf{Disaster Management} \\
\textbf{Organization:} & Ministry of Home Affairs \\
\textbf{Department:} & National Disaster Response Force (NDRF), DM Division \\
\textbf{Recommended Stack:} & \texttt{PyTorch, Scikit-Learn, XGBoost, Pandas, GeoPandas, PostGIS} \\
\textbf{Key Challenges:} & \textit{\small Data processing \& rule-based decision engine; Standard API feeds \& tabular time-series data; Cloud web dashboard \& RESTful API microservices; Geotechnical vulnerability \& multi-hazard spatial zoning} \\
\end{tabular}

\vspace{0.2cm}
\hrule
\vspace{0.2cm}

\textbf{\color{sihnavy}Problem Statement Scope \& Requirements:}\\[0.15cm]
\noindent\hspace*{1.2em}\hangindent=2.4em\hangafter=1\textcolor{sihblue}{$\bullet$}~\small Background India's disaster-prone regions face recurring hazards such as landslides,floods, coastal erosion, and cloudbursts. Vulnerable habitations often remain in unsafe zones, leading to repeated loss of lives and property.Current relocation efforts are largely reactive, initiated after disasters strike, rather than proactively planned.\par\vspace{0.05cm}
\noindent\hspace*{1.2em}\hangindent=2.4em\hangafter=1\textcolor{sihblue}{$\bullet$}~\small Description The initiative seeks to develop an intelligent, GIS-enabled decision support platform. This platform will dynamically identify and update multi-hazard Red Zones (areas unsuitable for permanent habitation),assess the carrying capacity of safer alternative sites, and prioritize vulnerable habitations for relocation. The system will integrate hazard intensity, population vulnerability, and disaster history to guide evidence-based decisions.\par\vspace{0.05cm}
\noindent\hspace*{1.2em}\hangindent=2.4em\hangafter=1\textcolor{sihblue}{$\bullet$}~\small Expected Solution A robust, AI-driven GIS platform that Maps and updates hazard-based Red Zones in real time, assesses suitability and carrying capacity of safer relocation sites, prioritizes vulnerable habitations for immediate,short-term, and medium-term relocation and provides actionable insights to State Disaster Management Authorities for proactive planning.\par\vspace{0.05cm}

\end{tcolorbox}
\vspace{0.4cm}
\needspace{6\baselineskip}
\phantomsection
\addcontentsline{toc}{subsection}{[5.6/10] [SIH26060] Digital Platform for efficient remote management of Indian Antarctic Resear...}

\begin{tcolorbox}[pscard, title={"\textbf{SIH26060} \quad $\vert$ \quad Digital Platform for efficient remote management of Indian Antarctic Resear...}]
\textbf{\large Digital Platform for efficient remote management of Indian Antarctic Research Stations}\\[0.25cm]

\begin{tabular}{@{}lp{12.5cm}@{}}
\textbf{Difficulty Rating:} & \textcolor{diffmod}{\textbf{\Large 5.6 / 10}} \quad (\textbf{Moderate (Intermediate)}) \\
\textbf{Problem ID:} & 26060 \quad $\vert$ \quad \textbf{Category:} \textbf{Software} \quad $\vert$ \quad \textbf{Theme:} \textbf{Disaster Management} \\
\textbf{Organization:} & Ministry of Earth Sciences (MoES) \\
\textbf{Department:} & National Centre for Polar andOcean Research (NCPOR) \\
\textbf{Recommended Stack:} & \texttt{PyTorch, Scikit-Learn, XGBoost, Pandas, GeoPandas, PostGIS} \\
\textbf{Key Challenges:} & \textit{\small Data processing \& rule-based decision engine; Standard API feeds \& tabular time-series data; Cloud web dashboard \& RESTful API microservices; Disaster response planning \& resource management} \\
\end{tabular}

\vspace{0.2cm}
\hrule
\vspace{0.2cm}

\textbf{\color{sihnavy}Problem Statement Scope \& Requirements:}\\[0.15cm]
\noindent\small Develop a Digital Twin framework for Maitri and Bharati stations integrating infrastructure, energy, logistics and environmental monitoring for efficient remote management.\par\vspace{0.06cm}

\end{tcolorbox}
\vspace{0.4cm}
\needspace{6\baselineskip}
\phantomsection
\addcontentsline{toc}{subsection}{[5.6/10] [SIH26074] Downscaling of weather forecast from Block level to Panchayat level: Inferr...}

\begin{tcolorbox}[pscard, title={"\textbf{SIH26074} \quad $\vert$ \quad Downscaling of weather forecast from Block level to Panchayat level: Inferr...}]
\textbf{\large Downscaling of weather forecast from Block level to Panchayat level: Inferring high-resolution plots/ data/ information from low-resolution plot /data /information /variables for agro-meteorological advisory services.}\\[0.25cm]

\begin{tabular}{@{}lp{12.5cm}@{}}
\textbf{Difficulty Rating:} & \textcolor{diffmod}{\textbf{\Large 5.6 / 10}} \quad (\textbf{Moderate (Intermediate)}) \\
\textbf{Problem ID:} & 26074 \quad $\vert$ \quad \textbf{Category:} \textbf{Software} \quad $\vert$ \quad \textbf{Theme:} \textbf{Disaster Management} \\
\textbf{Organization:} & Ministry of Earth Sciences (MoES) \\
\textbf{Department:} & India Meteorological Department \\
\textbf{Recommended Stack:} & \texttt{FastAPI / Django, PostgreSQL, React / Next.js} \\
\textbf{Key Challenges:} & \textit{\small Data processing \& rule-based decision engine; Standard API feeds \& tabular time-series data; Cloud web dashboard \& RESTful API microservices; Disaster response planning \& resource management} \\
\end{tabular}

\vspace{0.2cm}
\hrule
\vspace{0.2cm}

\textbf{\color{sihnavy}Problem Statement Scope \& Requirements:}\\[0.15cm]
\noindent\small Downscaling of weather forecast from Block level to Panchayat level: Inferring high-resolution plots/ data/ information from low-resolution plot /data /information /variables for agro-meteorological advisory services.\par\vspace{0.06cm}

\end{tcolorbox}
\vspace{0.4cm}
\needspace{6\baselineskip}
\phantomsection
\addcontentsline{toc}{subsection}{[5.6/10] [SIH26206] Student Innovation}

\begin{tcolorbox}[pscard, title={"\textbf{SIH26206} \quad $\vert$ \quad Student Innovation}]
\textbf{\large Student Innovation}\\[0.25cm]

\begin{tabular}{@{}lp{12.5cm}@{}}
\textbf{Difficulty Rating:} & \textcolor{diffmod}{\textbf{\Large 5.6 / 10}} \quad (\textbf{Moderate (Intermediate)}) \\
\textbf{Problem ID:} & 26206 \quad $\vert$ \quad \textbf{Category:} \textbf{Software} \quad $\vert$ \quad \textbf{Theme:} \textbf{Disaster Management} \\
\textbf{Organization:} & AICTE \\
\textbf{Department:} & AICTE, MIC-Student Innovation \\
\textbf{Recommended Stack:} & \texttt{FastAPI / Django, PostgreSQL, React / Next.js} \\
\textbf{Key Challenges:} & \textit{\small Data processing \& rule-based decision engine; Standard API feeds \& tabular time-series data; Cloud web dashboard \& RESTful API microservices; Disaster response planning \& resource management} \\
\end{tabular}

\vspace{0.2cm}
\hrule
\vspace{0.2cm}

\textbf{\color{sihnavy}Problem Statement Scope \& Requirements:}\\[0.15cm]
\noindent\small Disaster management includes ideas related to risk mitigation, Planning and management before, after or during a disaster.\par\vspace{0.06cm}

\end{tcolorbox}
\vspace{0.4cm}
\end{document}